\chapter{逻辑门与组合电路}

\begin{keyidea}
有了可靠电平，下一步是把受控路径组织成逻辑门，再把门组织成选择、译码和运算网络。本章始终同时追踪逻辑功能、传播延迟、负载和毛刺。
\end{keyidea}

\section{一、从路径表得到逻辑门}

门电路的第一步不是画符号，而是把输入组合列完整。两个输入 A、B 只有四种组合。电路必须对每一种组合给出明确输出，而且每一行都要满足前面的要求：输出节点不能悬空，不能上下冲突，必须能在规定时间内进入可靠电平区间。

从需求到门电路，建议固定采用五步。第一，给每个输入和输出写清楚硬件含义，特别注意别把低有效信号误读成高有效。第二，列出完整真值表，不跳过“看起来不会发生”的行——只按自然语言想象少数情况，是最常见的漏项来源。第三，标出输出为 1 的行，或者标出输出为 0 的禁止行，不遗漏边界输入和异常组合。第四，把每一行翻译成开关路径：哪些条件使上拉导通，哪些条件使下拉导通，不写只有公式、不看电流路径。第五，检查每一行是否有明确驱动、是否存在冲突、是否满足电平和延迟约束——只证明逻辑等价，不算证明硬件可靠。

以安全允许灯为例。A 表示保护盖合上，B 表示急停未按下。允许灯亮的条件很严格：任一条件不满足，输出都必须为 0。

\begin{longtable}{L{0.12\linewidth}L{0.18\linewidth}L{0.16\linewidth}L{0.42\linewidth}}
\toprule
A & B & 允许灯 & 为什么 \\
\midrule
0 & 0 & 0 & 盖子没合上，急停也不是可运行状态 \\
0 & 1 & 0 & 急停条件满足，但盖子没合上 \\
1 & 0 & 0 & 盖子合上了，但急停条件不满足 \\
1 & 1 & 1 & 两个安全条件同时满足 \\
\bottomrule
\end{longtable}

该表对应 AND 的行为。它不是抽象数学定义，而是“两个条件都满足才允许”的硬件约束。若某个实现让 A=1、B=0 时灯也亮，那就不是 AND 近似得不够好，而是安全行为错了。真值表的价值在于防止“只凭自然语言写电路”。“两个条件都满足”看起来简单，但实现时很容易漏掉某一行。

再看故障灯。A 表示温度报警，B 表示电流报警。故障灯的规则反过来：只要有一个坏消息出现，就要亮。

\begin{longtable}{L{0.12\linewidth}L{0.18\linewidth}L{0.16\linewidth}L{0.42\linewidth}}
\toprule
A & B & 故障灯 & 为什么 \\
\midrule
0 & 0 & 0 & 没有任何故障报警 \\
0 & 1 & 1 & 电流报警已经足够触发故障 \\
1 & 0 & 1 & 温度报警已经足够触发故障 \\
1 & 1 & 1 & 两个故障同时存在，仍应报警 \\
\bottomrule
\end{longtable}

该表对应 OR 的行为。AND 和 OR 的差别不在符号长什么样，而在需求不同：一个要求所有条件同时成立，一个要求任一条件成立即可。NOT 只有一个输入，但同样要列完整。若 A 表示“禁止信号”，输出 Y 表示“允许动作”，那么 A=0 时 Y=1，A=1 时 Y=0。它解决的是反向解释：输入出现时，输出反而关闭。

到这里六个基本门都已经出场。下面给出它们的标准符号：

\begin{circuit}[x=1cm,y=1cm]
  \node[not gate US, draw=BookAccent, line width=0.55pt, scale=1.35] (g1) at (0,1.8) {};
  \node[hw label] at (0.2,0.9) {NOT：Y = NOT A};
  \node[and gate US, draw=BookAccent, line width=0.55pt, logic gate inputs=nn, scale=1.35] (g2) at (3.4,1.8) {};
  \node[hw label] at (3.7,0.9) {AND：Y = A AND B};
  \node[or gate US, draw=BookAccent, line width=0.55pt, logic gate inputs=nn, scale=1.35] (g3) at (6.9,1.8) {};
  \node[hw label] at (7.2,0.9) {OR：Y = A OR B};
  \node[nand gate US, draw=BookAccent, line width=0.55pt, logic gate inputs=nn, scale=1.35] (g4) at (0,-0.5) {};
  \node[hw label] at (0.5,-1.4) {NAND：先 AND 再反相};
  \node[nor gate US, draw=BookAccent, line width=0.55pt, logic gate inputs=nn, scale=1.35] (g5) at (3.4,-0.5) {};
  \node[hw label] at (3.95,-1.4) {NOR：先 OR 再反相};
  \node[xor gate US, draw=BookAccent, line width=0.55pt, logic gate inputs=nn, scale=1.35] (g6) at (6.9,-0.5) {};
  \node[hw label] at (7.2,-1.4) {XOR：两输入不同为 1};
\end{circuit}
\diagramnote{六个基本门的标准符号（美式画法）。注意输出端的小圆圈：它表示反相——NAND 就是 AND 加一个小圆圈，NOR 就是 OR 加一个小圆圈。这个小圆圈和晶体管开关模型里 PMOS 控制端的圆圈是同一个约定：有圈即取反。}

\begin{warningbox}
不要把“0 是假、1 是真”机械套到每个硬件信号上。某些信号是低有效，例如 \code{reset_n}、\code{chip_select_n}，低电平才表示动作发生。读这类信号时，先写真值表，再谈门电路。
\end{warningbox}

低有效信号值得单独处理。名字中的 \code{\_n} 常表示 active low，即低电平有效。若信号名为 \code{emergency_ok_n}，那么 0 可能表示“急停链路正常允许运行”，1 反而表示“不允许”。这种命名并不适合作为入门案例的默认读法，但真实硬件中很常见，因为开漏输出、有线汇总、复位链路和片选信号都可能倾向于低有效形式。读低有效信号时，不能先把 1 写成“发生”、0 写成“未发生”；必须先读信号约定。

\begin{longtable}{L{0.22\linewidth}L{0.28\linewidth}L{0.4\linewidth}}
\toprule
信号名 & 有效电平 & 读法 \\
\midrule
\code{reset_n} & 0 & 低电平时复位动作发生 \\
\code{chip_select_n} & 0 & 低电平时芯片被选中 \\
\code{output_enable_n} & 0 & 低电平时输出驱动器打开 \\
\code{fault_n} & 0 & 低电平可能表示故障请求有效，需看约定 \\
\bottomrule
\end{longtable}

低有效并不只是命名习惯，它会改变推理方向。假设外部急停链路使用低有效输入 \code{emergency_stop_n}：链路正常释放时为 1，急停动作、断线或安全继电器打开时为 0。若内部逻辑需要的是正逻辑 \code{emergency_ok}，就必须先做一次语义转换：

\begin{lstlisting}
emergency_ok = emergency_stop_n
stop_request = NOT emergency_stop_n
\end{lstlisting}

这里 \code{emergency_ok} 和 \code{stop_request} 都可以是正确内部信号，但它们表达的事实不同。前者适合放进 \code{allow_start} 的 AND 链；后者适合放进故障或停机请求的 OR 链。若在同一个表达式里混用低有效原始信号和正逻辑语义信号，很容易写出方向相反的门网络。正式硬件文档通常会在信号表里明确有效电平，就是为了防止这种错误。

真值表说清楚“应该怎样”，路径表说明“为什么会这样”。用受控开关看 AND，最自然的结构是串联。两个开关串在同一条路径上，只有 A 和 B 都让自己的开关闭合时，整条路径才通。OR 更像并联。A 控制一条路径，B 控制另一条路径。只要任意一路闭合，输出就能被拉向目标电平。

\begin{longtable}{L{0.12\linewidth}L{0.12\linewidth}L{0.34\linewidth}L{0.32\linewidth}}
\toprule
A & B & 串联路径状态 & AND 输出为什么成立或不成立 \\
\midrule
0 & 0 & 两个开关都断开 & 没有完整导通路径 \\
0 & 1 & A 断开，B 闭合 & 串联路径仍被 A 切断 \\
1 & 0 & A 闭合，B 断开 & 串联路径仍被 B 切断 \\
1 & 1 & 两个开关都闭合 & 从源到输出的路径完整 \\
\bottomrule
\end{longtable}

如果只画一条“让输出为 1 的路径”，还没有完成门电路设计。输出节点不能靠悬空表示 0。因此每一行真值表还要补一张路径表，说明高电平路径和低电平路径的状态。一个合法门在稳定输入下，不应该让输出悬空，也不应该让上拉和下拉强路径同时导通。

从开关网络看，NAND 和 NOR 往往比 AND 和 OR 更像基础积木。原因是它们天然带反相输出，更容易形成互补的上拉和下拉结构：当下拉网络导通时输出被拉低；当下拉网络不导通时，上拉网络负责把输出拉高。

\begin{longtable}{L{0.1\linewidth}L{0.1\linewidth}L{0.22\linewidth}L{0.46\linewidth}}
\toprule
A & B & NAND 输出 & 开关直觉 \\
\midrule
0 & 0 & 1 & 下拉串联路径断开，上拉负责拉高 \\
0 & 1 & 1 & A 断开，下拉路径不完整 \\
1 & 0 & 1 & B 断开，下拉路径不完整 \\
1 & 1 & 0 & A、B 都闭合，下拉路径完整，输出被拉低 \\
\bottomrule
\end{longtable}

NOR 可以用同样方式读。NOR 的输出只有在 A=0 且 B=0 时为 1；只要 A 或 B 为 1，输出就为 0。若以下拉网络看，A、B 并联，只要任一路输入为 1，就能把输出拉低；只有两路都不导通时，上拉网络把输出拉高。

NAND 和 NOR 还可以帮助理解 De Morgan 定律。逻辑上，\code{NOT (A OR B)} 等价于 \code{(NOT A) AND (NOT B)}；硬件上，这句话的含义是：只要 A 或 B 有一路故障报警，输出就应被拉低；只有 A 和 B 都没有故障报警，输出才允许被拉高。对于安全启动控制板，若 \code{fault = temp_alarm OR current_alarm}，那么“没有故障”就是 \code{NOT fault}，也就是 \code{(NOT temp_alarm) AND (NOT current_alarm)}。这不是代数游戏，而是把“任一故障禁止启动”和“所有故障都不存在才允许启动”写成同一硬件约束的两种形式。

\begin{lstlisting}
fault = temp_alarm OR current_alarm
no_fault = NOT fault
no_fault = (NOT temp_alarm) AND (NOT current_alarm)
\end{lstlisting}

只用 NAND 就能构造 NOT、AND、OR：

\begin{lstlisting}
NOT A = A NAND A
A AND B = (A NAND B) NAND (A NAND B)
A OR B = (A NAND A) NAND (B NAND B)
\end{lstlisting}

这说明 NAND 在功能上完备，不说明所有实现都应该只用 NAND。真实设计会在功能正确的前提下，考虑门级数量、路径延迟、面积、功耗和负载。逻辑式等价，只说明稳定输出行为相同；它不保证延迟、功耗和驱动能力相同。

例如安全允许信号可以先写成：

\begin{lstlisting}
allow_start = start_request AND cover_closed
              AND emergency_ok AND no_fault
\end{lstlisting}

若只用 NAND 实现，需要把多输入 AND 拆成若干 NAND 与反相结构。每拆一次，逻辑功能仍可保持，但门级层数、负载和延迟都会变化。因此“用 NAND 可以实现”只是功能完备性结论；真正进入硬件设计时，还要继续检查最长路径和每一级扇出。

XOR 的输出在两个输入不同的时候为 1，相同的时候为 0。它解决的是“检测两个 bit 是否不同”这个具体问题。

\begin{longtable}{L{0.16\linewidth}L{0.16\linewidth}L{0.24\linewidth}L{0.34\linewidth}}
\toprule
A & B & A XOR B & 解释 \\
\midrule
0 & 0 & 0 & 两个输入相同 \\
0 & 1 & 1 & 两个输入不同 \\
1 & 0 & 1 & 两个输入不同 \\
1 & 1 & 0 & 两个输入相同 \\
\bottomrule
\end{longtable}

可以把 XOR 拆成两种会输出 1 的情况：A=1 且 B=0，或者 A=0 且 B=1：

\begin{lstlisting}
A XOR B = (A AND NOT B) OR (NOT A AND B)
\end{lstlisting}

这条式子不要只当代数记忆。它对应两个乘积项：一项在“只有 A 为 1”时成立，另一项在“只有 B 为 1”时成立，最后的 OR 把两项合并。若两个输入都为 0，两项都不成立；若两个都为 1，两个“只有一个为 1”的项同样都不成立。XOR 后面会反复出现：两个 bit 相加时，和位就是“两个输入是否不同”；比较两个 bit 时，XOR 可以告诉我们它们是否不同；接一个 NOT，就能得到“两个输入是否相同”。

XOR 也说明了“常用门”和“便宜门”不一定一致。在纸面表达式里，XOR 是一个符号；在门级实现里，它往往比 AND、OR、NOT 更复杂，路径延迟也可能更长。因此，读到加法器、比较器、校验位或地址译码时，不能只数逻辑符号数量。若关键路径里连续经过多个 XOR，速度代价可能比同样数量的简单门更高。经典教材在讲门延迟时强调这一点，是为了让读者把布尔函数和实现代价同时放进脑中。

\topic{XOR 与 XNOR：门级构成与三个经典用途}
XOR 既然比基本门贵，就值得先看清它贵在哪里。若器件库只有二输入 NAND，XOR 最少可以用 4 个 NAND 实现。办法不是凭空想的，而是把与或式逐步改写成 NAND 形式：

\begin{lstlisting}
n = A NAND B
p = A NAND n       -- 等价于 (NOT A) OR B
q = B NAND n       -- 等价于 A OR (NOT B)
Y = p NAND q       -- 展开后正是 (A AND NOT B) OR (NOT A AND B)
\end{lstlisting}

按信号流向数级数：\code{n} 在第一级，\code{p}、\code{q} 在第二级，\code{Y} 在第三级——一个二输入 XOR 需要三级门延迟。若改用与或式直接搭建，需要 2 个反相器、2 个 AND、1 个 OR 共 5 个门，信号同样要依次经过反相、AND、OR 三级。两种搭法殊途同归：XOR 无论怎么实现，门数和级数都比 AND、OR、NAND 多。XNOR 只需在 XOR 之后再取反一次，门级成本再加一级反相。这就是前文“常用门不等于便宜门”的具体含义：表达式里一个清爽的符号，在门级就是三级延迟。

值得付出这个代价，是因为 XOR 回答的是算术与校验中最常问的问题：“两个 bit 是否不同”。三个经典用途几乎贯穿本书后续所有章节。

第一个用途是奇偶校验。8 位数据的偶校验位就是把 8 个 bit 全部异或起来：1 的个数为偶数时结果为 0，为奇数时结果为 1。8 个输入两两分组，用 7 个二输入 XOR 排成 3 层树（第一层 4 个，第二层 2 个，第三层 1 个）；若图省事排成一条链，则要 7 级延迟。树形结构门数不变、级数从 7 降到 3，这是“同一函数、不同拓扑、不同延迟”的又一个实例。接收端把收到的 8 位再算一次同一棵树，与随数据传来的校验位比较，任何单个 bit 出错都会让两边不一致——XOR 树因此是通信和存储系统里最小的检错硬件。

第二个用途是加法器的和位。后文全加器的 \code{S = A XOR B XOR Cin}，问的正是“A、B、Cin 中是否有奇数个 1”；而进位 \code{Cout} 问的是“是否至少两个 1”，即三传感器判断器。一位全加器因此可以读成一个奇偶网络加一个多数网络：XOR 负责“奇数个”，AND/OR 负责“至少两个”。

第三个用途是相等比较。XOR 输出 1 表示“不同”，把它反过来：XNOR 输出 1 表示“相同”。比较两个 4 位数是否相等，就是逐位 XNOR，再把四路结果用 AND 归并：

\begin{lstlisting}
eq0 = A0 XNOR B0
eq1 = A1 XNOR B1
eq2 = A2 XNOR B2
eq3 = A3 XNOR B3
equal = eq0 AND eq1 AND eq2 AND eq3
\end{lstlisting}

归并的 AND 同样可以排成两级树而不是四输入单门，以控制扇入。后文数值比较器还会看到：相等信号 \code{eq_i} 不只用于判断全等，还是大小比较的裁决基础。

\begin{longtable}{L{0.14\linewidth}L{0.28\linewidth}L{0.3\linewidth}L{0.2\linewidth}}
\toprule
用途 & 回答的问题 & 关键结构 & 检查点 \\
\midrule
奇偶校验 & 1 的个数是奇还是偶 & n bit 用 n$-$1 个 XOR 排成 $\lceil\log_2 n\rceil$ 级树 & 排成链则延迟随位数线性增长 \\
加法器和位 & 是否有奇数个 1 & 两级 XOR：\code{A XOR B} 再 XOR \code{Cin} & 和位与进位是两条不同延迟的路径 \\
相等比较 & 每一位是否都相同 & XNOR 逐位加 AND 归并树 & 归并树分级，避免单门扇入过大 \\
\bottomrule
\end{longtable}

这张表也预告了一个反复使用的读图方法：看到 XOR，先问它在数“不同”还是数“奇偶”；看到 XOR 树，先数层数而不是门数。

真值表还可以用“最小项”来读。所谓最小项，就是与一行输入组合对应的完整乘积项。例如三输入 A、B、C 中，只有 A=1、B=0、C=1 这一行输出为 1，则对应的乘积项是 \code{A AND (NOT B) AND C}。把所有输出为 1 的行用 OR 合并，就得到一种直接实现。它不一定最简，但它有一个优点：每一项都能追溯到真值表中的具体一行。

\begin{longtable}{L{0.16\linewidth}L{0.2\linewidth}L{0.48\linewidth}}
\toprule
A B C & 输出 Y & 对应乘积项 \\
\midrule
0 1 1 & 1 & \code{(NOT A) AND B AND C} \\
1 0 1 & 1 & \code{A AND (NOT B) AND C} \\
1 1 0 & 1 & \code{A AND B AND (NOT C)} \\
1 1 1 & 1 & \code{A AND B AND C} \\
\bottomrule
\end{longtable}

若把这些最小项直接相加，就能实现“三个输入至少两个为 1”的判断。但前文给出的表达式 \code{(A AND B) OR (A AND C) OR (B AND C)} 更短，因为它把若干行合并成“任意两项同时成立”的条件。这里体现出两个层次：最小项保证从真值表机械得到正确表达式，化简则减少门数和路径深度。化简以后仍要确认没有改变低有效含义、默认安全策略和毛刺风险。

卡诺图的核心可以用“相邻最小项”的文字规则理解。若两个最小项只在一个输入上不同，且其余输入完全相同，那么这个不同输入对输出为 1 并不是必要条件，可以被消去。例如：

\begin{lstlisting}
(A AND B AND C) OR (A AND B AND NOT C) = A AND B
\end{lstlisting}

左边两项分别覆盖 C=1 和 C=0 两行；既然 C 的两种取值都让输出为 1，最终表达式就不必再依赖 C。卡诺图的格子相邻关系，本质上就是帮助设计者系统地寻找这种“只差一个输入”的合并机会。合并得当可以减少门数、缩短路径和降低负载；但合并之后必须重新检查低有效语义、非法状态和过渡行为。

文字规则必须落实到格子才算学会。以“三个输入至少两个为 1”为例，把 A 放在行、BC 放在列，每个格子填一行的输出：

\begin{circuit}[x=1cm,y=1cm]
  \node[hw label] at (-0.7,0.35) {A};
  \node[hw label] at (2.2,1.15) {BC};
  \foreach \x/\v in {0/00, 1.1/01, 2.2/11, 3.3/10} {
    \node[hw label] at (\x+0.55,0.35) {\v};
    \draw[BookMuted, line width=0.4pt] (\x,0) -- (\x,-2.2);
  }
  \node[hw label] at (4.95,0.35) { };
  \foreach \y/\v in {0/0, -1.1/1} {
    \node[hw label] at (-0.7,\y-0.55) {\v};
  }
  \draw[BookMuted, line width=0.4pt] (4.4,0) -- (4.4,-2.2);
  \draw[BookMuted, line width=0.4pt] (0,0) -- (4.4,0);
  \draw[BookMuted, line width=0.4pt] (0,-1.1) -- (4.4,-1.1);
  \draw[BookMuted, line width=0.4pt] (0,-2.2) -- (4.4,-2.2);
  \node at (0.55,-0.55) {0};
  \node at (1.65,-0.55) {0};
  \node at (2.75,-0.55) {1};
  \node at (3.85,-0.55) {0};
  \node at (0.55,-1.65) {0};
  \node at (1.65,-1.65) {1};
  \node at (2.75,-1.65) {1};
  \node at (3.85,-1.65) {1};
  \draw[BookTeal, line width=0.9pt, rounded corners=5pt] (2.3,0.12) rectangle (3.3,-2.32);
  \draw[BookAccent, line width=0.9pt, rounded corners=5pt] (1.15,-1.2) rectangle (3.4,-2.38);
  \draw[BookAmber, line width=0.9pt, rounded corners=5pt] (2.2,-1.05) rectangle (4.45,-2.25);
  \node[hw label, text=BookTeal] at (2.8,0.75) {消去 A：BC};
  \node[hw label, text=BookAccent] at (5.35,-1.35) {消去 B：AC};
  \node[hw label, text=BookAmber] at (5.35,-2.05) {消去 C：AB};
\end{circuit}
\diagramnote{三变量卡诺图（注意列头顺序 00、01、11、10 是 Gray 码：相邻两列永远只差一个 bit，这样“相邻”才有消元意义）。三个圈各合并两个相邻的 1：每个圈里那个“两种取值都出现”的输入被消去，得到 \code{Y = AB OR AC OR BC}——与后文三传感器判断器用最小项硬推出的结果一致，但快得多。}

化简还可能改变毛刺形态。某些冗余项对稳定真值表没有贡献，却能在输入翻转时提供连续路径。后文会看到，\code{(A AND B) OR ((NOT A) AND C)} 在 B=1、C=1 且 A 翻转时可能短暂掉到 0；增加共识项 \code{B AND C} 不改变稳定真值表，却能减少这类静态毛刺风险。因此，安全相关组合逻辑不应把“项数最少”当作唯一目标。形式上最简的表达式，未必是危险输出最稳妥的实现。

有时真值表里会出现“无关项”。所谓无关，不是设计者懒得定义，而是某些输入组合在系统约定中不会被使用，或者该组合出现时输出 0/1 都不会影响可观察行为。例如一个 2-bit 故障编码只允许 00、01、10，保留 11 作为非法状态；若后级在 11 时一定进入停机诊断，那么某个普通指示灯在 11 下亮或灭可能都可以接受。化简逻辑时可以利用无关项减少门数，但必须谨慎：一旦所谓“不可能”的输入来自外部连接器、未初始化状态或故障状态，它就不应轻易被当作无关项。

\begin{longtable}{L{0.2\linewidth}L{0.34\linewidth}L{0.36\linewidth}}
\toprule
输入组合类别 & 可以怎样处理 & 必须确认的问题 \\
\midrule
正常组合 & 严格定义输出 & 是否覆盖了所有业务行为 \\
保留组合 & 可定义为诊断、停机或无关项 & 后级是否真的不会误用该输出 \\
非法组合 & 通常导向安全默认值 & 硬件上是否可能短暂出现 \\
外部故障组合 & 不应随意当作无关项 & 断线、噪声、上电瞬间如何表现 \\
\bottomrule
\end{longtable}

对于安全启动控制板，保守做法是：任何未确认安全的组合，都不应让 \code{allow_start} 为 1。即使某个组合在正常操作流程中不会出现，只要它可能由断线、上电初态、输入调理失败或编码错误产生，就应进入“不启动”或“故障”一侧。组合逻辑的化简必须服从系统约定，而不是反过来让系统约定迁就最短表达式。

逻辑化简还要保留可审查性。对教学和安全控制而言，表达式稍长但能清楚对应需求，往往比极度压缩的表达式更容易审查。若某个输出必须由四个安全条件共同允许，那么把表达式保留为四个语义项，有助于后续检查每个条件来自哪里、默认值是什么、是否经过输入调理。化简可以减少门数和延迟，但不能让规格意图变得不可见。

四种常见表达方式各有取舍。逐项语义表达式容易对应需求和安全审查，但门数和层数通常不是最少；最小项展开可以机械追溯到真值表行，表达式却可能很长、延迟不一定好；代数化简式可能减少门数和路径层数，却可能隐藏低有效、默认值和故障策略；器件库映射式最接近实际可实现结构，但必须重新检查负载、扇入和时序。选择哪一种，取决于这段逻辑更需要可审查性还是更少的门级资源。

因此，正式设计常保留多个层次的描述：需求层写“所有安全条件满足才允许启动”；逻辑层写 \code{allow_start = start_request AND cover_closed AND emergency_ok AND no_fault}；实现层再说明由哪些门、缓冲和输出驱动器组成；验证层记录电平、延迟、毛刺和异常输入的检查结果。四个层次互相对应，才能避免“公式正确但系统含义错误”。

处理器控制逻辑中常见的 PLA，可以提前理解为“可规则化实现的组合逻辑”。PLA 的基本思想是先用一组 AND 项识别输入条件，再用 OR 项合成输出控制信号。例如指令译码可以把操作码、状态位和时序阶段转成若干控制线：选择哪一路进入 ALU、是否写寄存器、是否读内存、下一步进入哪个控制状态。数字硬件基础部分不需要展开 PLA 版图，但读者应看出它和最小项方法的关系：先识别条件，再汇总结果。

\begin{longtable}{L{0.2\linewidth}L{0.34\linewidth}L{0.36\linewidth}}
\toprule
PLA 视角 & 本章对应概念 & 在处理器中的用途 \\
\midrule
输入条件 & 真值表行、最小项、状态位 & 操作码、标志位、时序阶段 \\
AND 平面 & 乘积项 & 某条指令或某类条件成立 \\
OR 平面 & 乘积项汇总 & 产生写寄存器、读内存、选择 ALU 等控制线 \\
输出控制 & 组合逻辑输出 & 驱动数据通路选择器和使能信号 \\
\bottomrule
\end{longtable}

8086 这类早期微处理器内部就包含大量译码和控制逻辑。它不是一堆孤立门的随意堆叠，而是把操作码解释、微操作顺序、总线周期和算术逻辑组合成可重复执行的机器动作。处理器内部那些看似高级的控制信号，本质上也必须由可靠电平、门网络、组合路径和时序边界支撑。

\begin{chipcase}
8086 常被放在“完整 CPU”章节讲，但在数字硬件基础部分也能提前借用它的局部结构。它把取指相关的总线接口工作和执行相关的内部运算工作分开，并通过指令预取队列尝试让取指与执行重叠。这个案例说明：低延迟并不只来自某个门更快，也来自\hwtiming{把等待路径重叠}。不过预取队列不能消除所有等待；分支改变控制流、总线访问受阻或指令消费速度超过预取速度时，执行仍会遇到空洞。
\end{chipcase}

80386 的案例则说明另一件事：晶体管增加以后，芯片不只是把原有电路做得更快，还会把更多系统功能搬到片内。32-bit 数据通路、保护机制、地址转换和更复杂控制逻辑，都意味着更多组合判断、更多寄存器边界和更复杂的关键路径管理。处理器越复杂，越不能把“延迟”只理解成单个门的传播时间；控制路径、数据路径、存储路径和外部接口路径都要分别分析。

下面列出几个经典芯片的结构特征，把“门网络、组合控制、寄存器边界和片上集成”放到真实产品谱系中：

\begin{longtable}{L{0.16\linewidth}L{0.28\linewidth}L{0.46\linewidth}}
\toprule
芯片 & 先观察的结构特征 & 与本章概念的关系 \\
\midrule
MOS 6502 & 早期微处理器，外部总线与内部控制紧密配合 & 指令执行依赖译码、数据通路选择和控制时序 \\
Zilog Z80 & 经典 8-bit CPU，寄存器组和总线控制较丰富 & 组合译码与状态机共同组织内存和 I/O 访问 \\
Intel 8051 & 单片机把 CPU、存储器、I/O 和定时器放到同一芯片 & 片上集成减少板级连接，使控制信号更局部 \\
Intel 80486 & 在处理器内加入更强流水线、片上 cache 和浮点方向的集成 & 常用路径靠近执行核心，外部等待被部分削减 \\
Pentium & 双流水线、分支预测和更复杂执行控制成为重点 & 低延迟不仅靠门快，还靠预测、并行和路径重叠 \\
\bottomrule
\end{longtable}

\section{二、门网络构成组合逻辑}

组合逻辑的特点很直接：输出只由当前输入决定，不保存过去。它像一个硬件函数，但这个函数不是一步一步执行的过程，而是一张由门和连线组成的网络。输入变化后，信号沿多条路径同时传播；等延迟过去，输出稳定到当前输入对应的值。

设计组合逻辑时，可以把前一节的五步具体化为一条工作顺序。先列出输入输出，让每个 0/1 的含义明确；再列出行为，用真值表或条件表覆盖所有输入组合；然后写出逻辑表达式，让输出为 1 或 0 的每个条件都对应一条可实现路径；接着合成网络，用基本门、选择器和译码器明确每个输出的驱动者和选择者；最后检查物理限制——延迟、扇出、毛刺、电平余量，并说明后级在什么时候使用这个输出。这个顺序不要求一开始就写出最简表达式，而是先把需求完整展开，再逐步收缩为硬件网络。

先做一个三传感器判断器。输入 A、B、C 表示三个传感器，要求至少两个传感器为 1 时输出 1。这个需求不能只凭直觉画线，先列出行为：

\begin{longtable}{L{0.1\linewidth}L{0.1\linewidth}L{0.1\linewidth}L{0.16\linewidth}L{0.44\linewidth}}
\toprule
A & B & C & Y & 原因 \\
\midrule
0 & 0 & 0 & 0 & 一个条件都没有 \\
1 & 0 & 0 & 0 & 只有一个条件 \\
0 & 1 & 0 & 0 & 只有一个条件 \\
0 & 0 & 1 & 0 & 只有一个条件 \\
1 & 1 & 0 & 1 & A 和 B 同时满足 \\
1 & 0 & 1 & 1 & A 和 C 同时满足 \\
0 & 1 & 1 & 1 & B 和 C 同时满足 \\
1 & 1 & 1 & 1 & 三个都满足 \\
\bottomrule
\end{longtable}

从表中可以直接得到一个门网络：

\begin{lstlisting}
Y = (A AND B) OR (A AND C) OR (B AND C)
\end{lstlisting}

这条表达式包含三个乘积项。第一项在 A、B 同时为 1 时成立；第二项在 A、C 同时为 1 时成立；第三项在 B、C 同时为 1 时成立。最后的 OR 不是任意合并，而是在检查三个条件中是否至少有一个成立。按这种方式阅读，门网络就不只是公式变形，而是把真值表中每一种输出为 1 的原因接成硬件。

两个 bit 相加是第二个基本组合逻辑例子。输出不止一个 bit，因为 1+1 得到二进制 10，需要一个和位 S 和一个进位 C。

\begin{longtable}{L{0.16\linewidth}L{0.16\linewidth}L{0.22\linewidth}L{0.22\linewidth}}
\toprule
A & B & S & C \\
\midrule
0 & 0 & 0 & 0 \\
0 & 1 & 1 & 0 \\
1 & 0 & 1 & 0 \\
1 & 1 & 0 & 1 \\
\bottomrule
\end{longtable}

观察表格可以看到，S 在 A、B 不同时为 1，C 在 A、B 同时为 1：

\begin{lstlisting}
S = A XOR B
C = A AND B
\end{lstlisting}

这就是半加器。它已经完成了一位相加，但它还缺一个输入：来自低位的进位。只要要做多位加法，最低位之外的每一位都不能只看 A 和 B，还要看低一位是否进上来。

全加器输入为 A、B、Cin，输出为 S、Cout。可以先把 A 和 B 用半加器相加，得到临时和 S1 和进位 C1；再把 S1 和 Cin 用第二个半加器相加，得到最终和 S 和第二个进位 C2；最后把两个进位合并。

\begin{lstlisting}
S1 = A XOR B
C1 = A AND B
S  = S1 XOR Cin
C2 = S1 AND Cin
Cout = C1 OR C2
\end{lstlisting}

先别急着看电路，把三输入八行真值表列全。加法本质是“数 1 的个数”：一个 1 都没有，和为 0；恰好一个 1，和为 1 不进位；两个 1，和位变 0、产生进位；三个 1，和位为 1 且进位：

\begin{longtable}{L{0.09\linewidth}L{0.09\linewidth}L{0.11\linewidth}L{0.09\linewidth}L{0.11\linewidth}L{0.39\linewidth}}
\toprule
A & B & \code{Cin} & S & \code{Cout} & 解释 \\
\midrule
0 & 0 & 0 & 0 & 0 & 没有 1 \\
0 & 0 & 1 & 1 & 0 & 只有一个 1（来自进位输入） \\
0 & 1 & 0 & 1 & 0 & 只有一个 1 \\
0 & 1 & 1 & 0 & 1 & 两个 1：本位写 0，向高位进 1 \\
1 & 0 & 0 & 1 & 0 & 只有一个 1 \\
1 & 0 & 1 & 0 & 1 & 两个 1 \\
1 & 1 & 0 & 0 & 1 & 两个 1 \\
1 & 1 & 1 & 1 & 1 & 三个 1：和位 1，再进 1 \\
\bottomrule
\end{longtable}

对照表格验证上面的式子：S 在“恰好一个 1”或“三个 1”时为 1，正是 \code{A XOR B XOR Cin}；\code{Cout} 在“至少两个 1”时为 1，正是三传感器判断器的结果。再把式子落成门，就是两个半加器加一个 OR：

\begin{circuit}[x=1cm,y=1cm]
  \node[xor gate US, draw=BookAccent, line width=0.5pt, logic gate inputs=nn, scale=1.2] (x1) at (2.0,2.9) {};
  \node[and gate US, draw=BookAccent, line width=0.5pt, logic gate inputs=nn, scale=1.2] (a1) at (2.0,1.2) {};
  \node[xor gate US, draw=BookAccent, line width=0.5pt, logic gate inputs=nn, scale=1.2] (x2) at (4.9,2.6) {};
  \node[and gate US, draw=BookAccent, line width=0.5pt, logic gate inputs=nn, scale=1.2] (a2) at (4.9,0.7) {};
  \node[or gate US, draw=BookAccent, line width=0.5pt, logic gate inputs=nn, scale=1.2] (o1) at (7.4,1.65) {};
  \draw[BookMuted, dashed, line width=0.45pt, rounded corners=4pt] (1.05,3.6) rectangle (2.95,0.5);
  \node[hw label] at (2.0,3.85) {半加器 1};
  \draw[BookMuted, dashed, line width=0.45pt, rounded corners=4pt] (3.95,3.3) rectangle (5.85,0.0);
  \node[hw label] at (4.9,3.55) {半加器 2};
  \node[hw label] at (-0.6,3.05) {A};
  \node[hw label] at (-0.6,2.35) {B};
  \node[hw label] at (-0.6,0.2) {\code{Cin}};
  \draw[hw bus] (-0.3,3.05) -- (x1.input 1);
  \draw[hw bus] (-0.3,2.35) -- (x1.input 2);
  \draw[fill=BookInk] (0.5,3.05) circle (0.05);
  \draw[hw bus] (0.5,3.05) -- (0.5,1.4) -- (a1.input 1);
  \draw[fill=BookInk] (0.8,2.35) circle (0.05);
  \draw[hw bus] (0.8,2.35) -- (0.8,1.0) -- (a1.input 2);
  \draw[hw bus] (x1.output) -- node[hw label, above]{\code{S1}} (x2.input 1);
  \draw[fill=BookInk] (3.3,2.9) circle (0.05);
  \draw[hw bus] (3.3,2.9) -- (3.3,0.9) -- (a2.input 1);
  \draw[hw bus] (-0.3,0.2) -- (3.6,0.2) -- (3.6,2.3) -- (x2.input 2);
  \draw[fill=BookInk] (3.6,0.5) circle (0.05);
  \draw[hw bus] (3.6,0.5) -- (a2.input 2);
  \draw[hw bus] (x2.output) -- node[hw label, above]{S} ++(1.3,0);
  \draw[hw bus] (a1.output) -- node[hw label, above]{\code{C1}} (6.4,1.2) -- (6.4,1.5) -- (o1.input 2);
  \draw[hw bus] (a2.output) -- node[hw label, above]{\code{C2}} (6.8,0.7) -- (6.8,1.8) -- (o1.input 1);
  \draw[hw bus] (o1.output) -- node[hw label, above]{\code{Cout}} ++(1.3,0);
\end{circuit}
\diagramnote{全加器 = 两个半加器 + 一个 OR。第一个半加器算 A+B，第二个半加器把临时和 S1 再与进位输入 Cin 相加；两处的进位输出经 OR 合并——只要有一处产生进位，就向高位进 1。把这张图的每个输出和上面八行真值表逐行对一遍，比背公式可靠得多。}

把全加器一位一位串起来，就得到串行进位加法器。它的逻辑很清楚：第 0 位算出进位，传给第 1 位；第 1 位再把进位传给第 2 位。问题也很清楚：若最低位产生的进位一路传到最高位，最高位必须等前面所有进位稳定后才能稳定。这是组合逻辑里第一次真正遇到“逻辑正确”和“速度足够”的差别。真值表保证最后结果正确，传播延迟决定结果多久之后才可靠。

\begin{circuit}[x=1cm,y=1cm]
  \foreach \x/\i in {0/0, 3/1, 6/2, 9/3} {
    \node[hw compute, minimum width=1.7cm, minimum height=1.5cm] (fa\i) at (\x,0) {FA\i\\全加器};
    \node[hw label] at (\x,1.35) {A\i, B\i};
    \draw[hw bus] (\x,1.05) -- (fa\i.north);
    \draw[hw bus] (fa\i.south) -- node[hw label, right]{S\i} (\x,-1.35);
  }
  \node[hw label] at (-1.6,0) {\code{Cin}=0};
  \draw[hw bus] (-1.1,0) -- (fa0.west);
  \draw[hw bus] (fa0.east) -- node[hw label, above]{\code{C1}} (fa1.west);
  \draw[hw bus] (fa1.east) -- node[hw label, above]{\code{C2}} (fa2.west);
  \draw[hw bus] (fa2.east) -- node[hw label, above]{\code{C3}} (fa3.west);
  \draw[hw bus] (fa3.east) -- node[hw label, above]{\code{Cout}} ++(1.1,0);
  \draw[hw danger] ($(fa0.south east)+(0.1,-0.5)$) -- node[hw label, text=BookRed, below, text width=6.5cm]{最长路径：最低位的进位要逐位穿过 4 级全加器，最高位才能稳定} ($(fa3.south east)+(0.1,-0.5)$);
\end{circuit}
\diagramnote{4 位串行进位加法器。每位的和位 S 只依赖本位输入，很快；但 C4 依赖 C3、C3 依赖 C2……进位像接力一样从低位向高位逐级传播，红色线段就是决定全机速度的那条最长路径。加法器位数越多，这条链越长——这正是数值与数据通路部分要解决的矛盾。}

全加器也可以用生成与传播来阅读。若 A 和 B 同时为 1，本位一定产生进位，称为生成；若 A 和 B 恰好一个为 1，那么本位不会自己产生进位，但会把输入进位传到输出，称为传播。

\begin{lstlisting}
generate  = A AND B
propagate = A XOR B
Cout = generate OR (propagate AND Cin)
S = propagate XOR Cin
\end{lstlisting}

这种写法为后面的加法器优化留下接口。串行进位把进位逐位传递，结构简单但最长路径随位数增长；更快的加法器会提前计算若干位的生成和传播关系，减少等待层数。组合电路不仅有“算什么”，还有“最慢路径从哪里穿过”。算术电路尤其如此，因为进位路径常常决定 ALU 的速度上限。

把全加器再向前推进一步，就得到 1-bit ALU 切片的雏形。ALU 不是不可分解的黑箱，而是许多位宽相同的切片并排工作：每一位接收本位的 A、B、来自低位的进位，以及操作选择信号；它同时准备若干候选结果，例如 AND、OR、XOR 和 ADD；最后由选择器决定哪一种候选结果成为本位输出（选择器是由控制信号从多路输入中选一路输出的组合电路，下一小节马上讲到）。进位输出则继续传给高一位或交给更快的进位逻辑处理。

\begin{lstlisting}
and_result = A AND B
or_result  = A OR B
xor_result = A XOR B

sum_result = A XOR B XOR Cin
Cout = (A AND B) OR ((A XOR B) AND Cin)

Y = mux(op, and_result, or_result, xor_result, sum_result)
\end{lstlisting}

这段行为式只是说明逻辑关系，不表示硬件按行执行。真实电路中，AND、OR、XOR 和加法候选路径可以同时传播；\code{op} 控制选择器，让其中一路成为输出。于是，一个 1-bit ALU 切片至少包含三类组合结构：基本逻辑门、全加器、选择器。多个切片并排以后，AND/OR/XOR 这类逐位运算的延迟主要取决于本位门和选择器；ADD 则还受进位路径影响。

\begin{circuit}[x=1cm,y=1cm]
  \node[hw compute, minimum width=1.5cm, minimum height=0.72cm] (and) at (1.2,2.7) {AND};
  \node[hw compute, minimum width=1.5cm, minimum height=0.72cm] (or) at (1.2,1.9) {OR};
  \node[hw compute, minimum width=1.5cm, minimum height=0.72cm] (xor) at (1.2,1.1) {XOR};
  \node[hw compute, minimum width=1.5cm, minimum height=1.0cm] (fa) at (1.2,0.0) {全加器\\ADD};
  \draw[line width=0.55pt] (4.2,3.3) -- (5.2,3.0) -- (5.2,-0.6) -- (4.2,-0.9) -- cycle;
  \node[hw label] at (4.7,3.6) {4 选 1 选择器};
  \node[hw label] at (-1.7,3.3) {A};
  \draw[hw bus] (-1.4,3.3) -- (-0.6,3.3) -- (-0.6,0.0);
  \node[hw label] at (-1.7,2.55) {B};
  \draw[hw bus] (-1.4,2.55) -- (-0.15,2.55) -- (-0.15,-0.3);
  \draw[fill=BookInk] (-0.6,2.7) circle (0.05);
  \draw[hw bus] (-0.6,2.7) -- (and.north west);
  \draw[fill=BookInk] (-0.15,2.7) circle (0.05);
  \draw[hw bus] (-0.15,2.7) -- (and.south west);
  \draw[fill=BookInk] (-0.6,1.9) circle (0.05);
  \draw[hw bus] (-0.6,1.9) -- (or.north west);
  \draw[fill=BookInk] (-0.15,1.9) circle (0.05);
  \draw[hw bus] (-0.15,1.9) -- (or.south west);
  \draw[fill=BookInk] (-0.6,1.1) circle (0.05);
  \draw[hw bus] (-0.6,1.1) -- (xor.north west);
  \draw[fill=BookInk] (-0.15,1.1) circle (0.05);
  \draw[hw bus] (-0.15,1.1) -- (xor.south west);
  \draw[hw bus] (-0.6,0.3) -- (fa.north west);
  \draw[fill=BookInk] (-0.6,0.3) circle (0.05);
  \draw[hw bus] (-0.15,-0.3) -- (fa.south west);
  \draw[fill=BookInk] (-0.15,-0.3) circle (0.05);
  \node[hw label] at (-1.7,-0.3) {\code{Cin}};
  \draw[hw bus] (and.east) -- (4.2,2.7);
  \draw[hw bus] (or.east) -- (4.2,1.9);
  \draw[hw bus] (xor.east) -- (4.2,1.1);
  \draw[hw bus] (fa.east) -- (4.2,0.0) -- (4.2,0.3);
  \draw[hw bus] (fa.south) -- node[hw label, right]{\code{Cout}：传给高一位切片} (1.2,-1.55);
  \node[hw label] at (4.7,-1.35) {\code{op}};
  \draw[hw bus] (4.7,-1.1) -- (4.7,-0.75);
  \draw[hw bus] (5.2,1.2) -- node[hw label, above]{Y} (6.3,1.2);
\end{circuit}
\diagramnote{1-bit ALU 切片：四路候选同时传播，\code{op} 控制选择器让其中一路成为本位输出 Y；进位不进选择器，沿进位链直接传给高位切片。把这种切片并排 16 个就是 8086 的加法/逻辑单元雏形，并排 32 个、64 个，就逼近 80386 和现代执行单元——结构没变，只是宽度和进位优化在变。}

减法通常不需要另造一套独立减法器。二进制补码中，\code{A - B} 可以写成 \code{A + (NOT B) + 1}。硬件上常见做法是在 B 输入前放置一组可控反相器或 XOR 门：做加法时让 B 原样进入全加器，做减法时让 B 翻转，同时把最低位进位输入置为 1。这样同一条加法进位链既可服务 ADD，也可服务 SUB。

\begin{lstlisting}
B_eff = B XOR B_invert
sum_result = A XOR B_eff XOR Cin
Cout = (A AND B_eff) OR ((A XOR B_eff) AND Cin)
\end{lstlisting}

\begin{longtable}{L{0.18\linewidth}L{0.2\linewidth}L{0.18\linewidth}L{0.34\linewidth}}
\toprule
操作 & \code{B_invert} & 最低位 \code{Cin} & 本位输出含义 \\
\midrule
ADD & 0 & 0 & \code{A + B} 的和位 \\
SUB & 1 & 1 & \code{A + (NOT B) + 1}，即补码减法 \\
AND & 不使用或置 0 & 不使用 & 输出 \code{A AND B}，进位链不作为结果来源 \\
OR & 不使用或置 0 & 不使用 & 输出 \code{A OR B}，进位链不作为结果来源 \\
XOR & 不使用或置 0 & 不使用 & 输出 \code{A XOR B}，进位链不作为结果来源 \\
\bottomrule
\end{longtable}

这张表说明\hwkey{控制信号本身也是硬件约定}。若 \code{op} 选择 SUB，但 \code{B_invert} 或最低位 \code{Cin} 没有同步改变，结果会变成错误的加法变体；若多位减法要生成借位、进位、零标志、符号标志或溢出标志，还要明确每个标志按哪一种数值解释计算。数字硬件基础部分不展开完整标志逻辑，只强调复用加法器并不意味着控制可以随意；控制位、数据路径和标志解释必须一起确定。

逐位逻辑操作路径最短：AND、OR、XOR 的本位输出只经过一两个门，主要受负载和选择器影响；它们适合逐位屏蔽、合并标志位、差异检测和翻转控制。加法路径则由全加器和进位链组成，最长路径常由进位传播决定。选择器一层看似透明，却是最容易被忽视的约定：\code{op} 必须稳定，不能让错误候选短暂输出到危险后级。这个切片模型提前解释了 CPU 数据通路里的三个常见事实。第一，位宽增加不只是“多写几个 bit”，而是复制更多切片，并重新处理进位、布线和负载。第二，控制信号不是附属注释；\code{op}、进位选择、结果写回选择都会直接改变哪条硬件路径有效。第三，ALU 的延迟不是单个门延迟，而是由候选路径、进位路径、选择器和输出负载共同决定。8086 的 16-bit 数据通路、80386 的 32-bit 数据通路，以及现代 64-bit 执行单元，都可以从这个切片观点继续放大理解。

多路选择器解决“多个输入选一个输出”。2 选 1 选择器的行为是：

\begin{lstlisting}
if sel = 0, Y = A
if sel = 1, Y = B
\end{lstlisting}

上述 \code{if} 只是描述行为，电路本身没有顺序执行。它的门级表达式是：

\begin{lstlisting}
Y = (NOT sel AND A) OR (sel AND B)
\end{lstlisting}

\begin{circuit}[x=1cm,y=1cm]
  \node[not gate US, draw=BookAccent, line width=0.5pt, scale=1.15] (n1) at (1.3,-1.0) {};
  \node[and gate US, draw=BookAccent, line width=0.5pt, logic gate inputs=nn, scale=1.2] (and1) at (2.9,2.3) {};
  \node[and gate US, draw=BookAccent, line width=0.5pt, logic gate inputs=nn, scale=1.2] (and2) at (2.9,0.1) {};
  \node[or gate US, draw=BookAccent, line width=0.5pt, logic gate inputs=nn, scale=1.2] (or1) at (5.3,1.2) {};
  \node[hw label] at (-1.1,2.6) {A};
  \draw[hw bus] (-0.8,2.6) -- (and1.input 1);
  \node[hw label] at (-1.1,0.4) {B};
  \draw[hw bus] (-0.8,0.4) -- (and2.input 1);
  \node[hw label] at (-1.1,-1.0) {\code{sel}};
  \draw[hw bus] (-0.8,-1.0) -- (n1.input);
  \draw[fill=BookInk] (0.2,-1.0) circle (0.05);
  \draw[hw bus] (0.2,-1.0) -- (0.2,-1.7) -- (1.7,-1.7) -- (and2.input 2);
  \draw[hw bus] (n1.output) -- (2.1,-1.0) -- (2.1,1.9) -- (and1.input 2);
  \draw[hw bus] (and1.output) -- (or1.input 1);
  \draw[hw bus] (and2.output) -- (or1.input 2);
  \draw[hw bus] (or1.output) -- node[hw label, above]{Y} (6.7,1.2);
  \node[hw label, text width=3.2cm] at (5.0,-0.9) {\code{sel}=0：上路开门，Y=A\\\code{sel}=1：下路开门，Y=B};
\end{circuit}
\diagramnote{2 选 1 选择器的门级实现：\code{sel} 直接控制下路 AND，反相后控制上路 AND——任何瞬间两路不会同时开门，输出永远有明确来源。选择器的“选择”不是程序分支，而是两扇由互补信号控制的门。}

选择器的作用很基础。若一条输出线可能来自加法器，也可能来自传感器，也可能来自某个寄存器，那么必须有一个明确的选择信号决定哪一路接到输出。没有选择器，多路来源同时驱动同一条线就会冲突。

选择器的控制信号本身也需要检查。若 \code{sel} 悬空，输出可能在 A 和 B 之间随机变化；若选择信号来自毛刺严重的组合网络，输出可能短暂选择错误来源；若 A、B 属于不同电源域或不同稳定时间，选择器输出还会继承这些边界问题。因此，选择器不是“安全地把两根线接在一起”，而是“在一个可靠控制信号约束下，只让一路值成为输出”。

\begin{longtable}{L{0.2\linewidth}L{0.34\linewidth}L{0.36\linewidth}}
\toprule
选择器检查项 & 问题 & 失败后果 \\
\midrule
选择信号默认值 & 上电或复位前选哪一路 & 输出来源不确定 \\
数据输入稳定性 & 被选中的输入是否已稳定 & 选择了正确来源但值仍不可靠 \\
未选输入行为 & 未选输入是否允许变化 & 不应通过内部耦合影响输出 \\
输出负载 & 选择器是否能驱动后级 & 边沿变慢或电平不足 \\
\bottomrule
\end{longtable}

译码器做相反方向的事：把一个二进制编号展开成 one-hot 信号（任意时刻恰好一条线为 1）。2-bit 输入可以选中 4 条输出中的一条。

\begin{longtable}{L{0.18\linewidth}L{0.48\linewidth}}
\toprule
输入 & 输出 \\
\midrule
00 & Y0=1，Y1/Y2/Y3=0 \\
01 & Y1=1，其他为 0 \\
10 & Y2=1，其他为 0 \\
11 & Y3=1，其他为 0 \\
\bottomrule
\end{longtable}

选择器把多路收成一路，译码器把编号展开成多路选择。

选择器还可以直接实现布尔函数。若一个函数有两个输入 A、B，可以用 A 作为选择信号，选择器的两个数据输入分别接“当 A=0 时 Y 应该等于什么”和“当 A=1 时 Y 应该等于什么”。例如 XOR 中，当 A=0 时，Y=B；当 A=1 时，Y=NOT B。因此一个 2 选 1 选择器加一个反相器即可实现 XOR：

\begin{lstlisting}
if A = 0: Y = B
if A = 1: Y = NOT B
Y = mux(A, B, NOT B)
\end{lstlisting}

ALU 输入选择、写回选择、下一 PC 选择，本质上都可以看成“控制信号选择哪一路值成为输出”。选择器不是附属小部件，而是把多种候选硬件动作收束成一个当前动作的基本结构。

\topic{多路复用器：树形扩展与任意函数实现}
一个 2 选 1 选择器只是一粒种子。回看它的门级结构：一个反相器、两个 AND、一个 OR——成本固定。把 2 选 1 当积木按树形堆叠，就能得到任意规模的选择器。4 选 1 用 3 个 2 选 1 排成两级：第一级两个，由选择信号低位 \code{S0} 分别在 \code{D0/D1}、\code{D2/D3} 中各选一路；第二级一个，由高位 \code{S1} 在两路胜者中选最终输出：

\begin{lstlisting}
w0 = mux(S0, D0, D1)
w1 = mux(S0, D2, D3)
Y  = mux(S1, w0, w1)
\end{lstlisting}

逐一代入可验证：\code{S1S0}=00 时 \code{Y=D0}，01 时 \code{Y=D1}，10 时 \code{Y=D2}，11 时 \code{Y=D3}。同一规律继续放大：8 选 1 用 7 个 2 选 1 排成三级，16 选 1 用 15 个排成四级。一般地，$2^n$ 选 1 需要 $2^n - 1$ 个 2 选 1，共 $n$ 级：

\begin{longtable}{L{0.14\linewidth}L{0.2\linewidth}L{0.14\linewidth}L{0.42\linewidth}}
\toprule
规模 & 2 选 1 个数 & 级数 & 结构读法 \\
\midrule
2 选 1 & 1 & 1 & 最小单元 \\
4 选 1 & 3 & 2 & 第一级两个并行，第二级一个汇总 \\
8 选 1 & 7 & 3 & 第一级四个并行，逐级减半 \\
16 选 1 & 15 & 4 & 数据输入指数增长，级数只线性增长 \\
\bottomrule
\end{longtable}

这张表藏着一个工程读法：选择器规模翻倍，数据输入数量翻倍，但延迟只增加一级——树形结构把指数增长的数据宽度换成了线性增长的级数。布线时还有一个讲究：最晚稳定的那一路选择信号应放在最后一级，让先到的数据先在低级门里传播。这与后文把晚到信号放在最后一级的做法是同一个原则。

多路复用器的第二个身份是通用函数发生器。一个 n 变量函数，可以把全部 n 个变量接到 $2^n$ 选 1 的选择端，再把真值表逐行抄到数据端：第 k 个数据端接第 k 行的输出值。这样实现函数不需要任何化简——选择器本身就是一张接死的真值表。更省的办法是只把 n$-$1 个变量接选择端，剩下的一个变量以四种形态之一出现在数据端：0、1、变量本身、变量的反。以三变量奇校验函数 \code{F(A,B,C)}（A、B、C 中有奇数个 1 时输出 1）为例，先用 8 选 1 实现，选择端 S2S1S0 接 A、B、C：

\begin{longtable}{L{0.14\linewidth}L{0.16\linewidth}L{0.12\linewidth}L{0.48\linewidth}}
\toprule
A B C & 最小项 & F & 数据端接法 \\
\midrule
0 0 0 & m0 & 0 & \code{D0} 接 0 \\
0 0 1 & m1 & 1 & \code{D1} 接 1 \\
0 1 0 & m2 & 1 & \code{D2} 接 1 \\
0 1 1 & m3 & 0 & \code{D3} 接 0 \\
1 0 0 & m4 & 1 & \code{D4} 接 1 \\
1 0 1 & m5 & 0 & \code{D5} 接 0 \\
1 1 0 & m6 & 0 & \code{D6} 接 0 \\
1 1 1 & m7 & 1 & \code{D7} 接 1 \\
\bottomrule
\end{longtable}

再省一级：只把 A、B 接 4 选 1 的选择端，数据端改成 C 的函数。逐行检查每个 A、B 组合下 F 随 C 的变化：

\begin{longtable}{L{0.14\linewidth}L{0.18\linewidth}L{0.18\linewidth}L{0.4\linewidth}}
\toprule
A B & C=0 时 F & C=1 时 F & 数据端接法 \\
\midrule
0 0 & 0 & 1 & \code{D0} 接 C \\
0 1 & 1 & 0 & \code{D1} 接 NOT C \\
1 0 & 1 & 0 & \code{D2} 接 NOT C \\
1 1 & 0 & 1 & \code{D3} 接 C \\
\bottomrule
\end{longtable}

例如 A=0、B=1 时：C=0 对应最小项 m2 输出 1，C=1 对应 m3 输出 0，F 恰好等于 NOT C，于是 \code{D1} 接反相后的 C——其余三行同理。一片 4 选 1 加一个反相器，就装下了整张八行真值表。

这个身份的意义远超节省门数。可编程逻辑器件中的查找表，本质就是把数据端换成可写存储单元的多路复用器：改写存储单元，就改写了函数。读法上也要注意代价的另一面：用选择器实现函数是机械的，它不利用函数内部结构——八行真值表和八个随机值，对 8 选 1 来说是同一个价格。化简的机会被换成了实现的整齐。

选择器用集中的逻辑裁决“谁成为输出”，前提是所有候选来源都汇聚到一处。当候选来源分散在不同芯片、不同板卡上时，还有另一种组织方式：让每个来源挂在同一条线上，各自决定何时驱动、何时松手。接下来的两个专题分别考察这种方式的两种输出级：三态与开漏。

\topic{三态输出与总线争用}
到目前为止，每个输出级都遵循同一约定：任何时刻要么强驱动高电平，要么强驱动低电平，绝不允许悬空，也绝不允许上下同时导通。三态输出在这个约定上增加第三个合法状态：高阻态。使能端无效时，输出级的上拉管和下拉管同时关断，输出端既不供出电流也不吸收电流，对外表现为“像没接一样”——它不是 0，也不是 1，而是退出。使能端有效时，它就是一个普通的推挽输出。按惯例，使能常做成低有效，记为 \code{OE_n}，与前文低有效信号约定一致：

\begin{longtable}{L{0.14\linewidth}L{0.12\linewidth}L{0.12\linewidth}L{0.52\linewidth}}
\toprule
\code{OE_n} & D & Y & 输出级状态 \\
\midrule
0 & 0 & 0 & 下拉管导通，强驱动低电平 \\
0 & 1 & 1 & 上拉管导通，强驱动高电平 \\
1 & 任意 & Z & 两管都关断，高阻，退出驱动 \\
\bottomrule
\end{longtable}

高阻态的价值在“共享”。若加法器、寄存器、外部接口都要把结果送到同一条数据线上，集中式选择器要求所有来源汇聚到一点；而给每个来源配一个三态输出，让它们平时处于高阻、轮到自己时才使能，同一条线就能被多个模块分时使用——这就是总线。总线能成立，靠的是一条硬约定：同一时刻最多一个驱动者。使能信号必须互斥，通常由译码器产生的 one-hot 或专门的仲裁逻辑保证。

违反这条约定的后果值得用电流算一遍。设两个三态驱动者同时使能，一个想推高、一个想拉低：电流从电源出发，经前者的上拉管、总线、后者的下拉管直到地，形成一条低阻通路。以每个输出级导通电阻约 25$\Omega$ 估算，这条通路上的电流约为 $5\mathrm{V} / (25\Omega + 25\Omega) = 100\mathrm{mA}$——比正常信号电流大一个数量级以上。此时总线电压停在中点附近，接收端读到的既不是合法高电平也不是合法低电平；大电流持续流过输出级，带来发热，反复或长时间争用会损坏器件。即使两个驱动者想驱动同一电平，超出额定的并联驱动也会使电平和边沿偏离数据手册的保证值。这就是总线争用：它不是逻辑错误——逻辑式里两个驱动者的值可以都正确；它是电气事故。

把这条冲突通路和正确的使能切换顺序画在一起，事故与预防各是什么模样就一目了然：

\begin{pdfdiagram}[x=0.9cm,y=0.9cm]
  % 上：两个三态驱动者对打的冲突电流通路（只画导通路径上的管子）
  \node[hw label] at (-0.3,6.2) {VCC};
  \draw[line width=0.55pt] (0.1,6.2) -- (4.5,6.2);
  \pic at (1.2,5.55) {hw pmos};
  \draw[line width=0.55pt] (1.2,5.97) -- (1.2,6.2);
  \pic at (3.4,4.08) {hw nmos};
  \pic at (3.4,3.2) {hw gnd};
  \draw[BookRed, line width=1.1pt] (1.2,5.13) -- (1.2,4.5) -- (3.4,4.5);
  \draw[BookRed, line width=1.1pt] (3.4,3.66) -- (3.4,3.2);
  \draw[hw danger] (1.2,5.0) -- (1.2,4.66);
  \draw[hw danger] (2.0,4.5) -- (2.36,4.5);
  \draw[hw danger] (3.4,3.6) -- (3.4,3.34);
  \node[hw label, anchor=west] at (1.75,5.55) {驱动器 A 上拉管导通（想推高）};
  \node[hw label, anchor=east] at (2.6,4.08) {驱动器 B 下拉管导通（想拉低）};
  \node[hw label] at (2.3,4.86) {总线};
  \node[hw label, text=BookRed, anchor=west] at (3.95,4.5) {冲突电流约 100mA};
  % 下：先断后通的使能切换时序（低有效使能）
  \draw[BookMuted, dashed, line width=0.45pt] (4,-1.35) -- (4,2.3);
  \draw[BookMuted, dashed, line width=0.45pt] (7,-1.35) -- (7,2.3);
  \draw[BookAccent, line width=0.9pt] (0,1.6) -- (4,1.6) -- (4,2.0) -- (12,2.0);
  \node[hw label, anchor=east] at (-0.2,1.8) {$\overline{OE_A}$};
  \draw[BookAccent, line width=0.9pt] (0,1.0) -- (7,1.0) -- (7,0.6) -- (12,0.6);
  \node[hw label, anchor=east] at (-0.2,0.8) {$\overline{OE_B}$};
  \draw[BookTeal, line width=0.9pt] (0,-1.0) rectangle (4,-0.3);
  \node at (2,-0.65) {A 驱动};
  \draw[BookMuted, line width=0.65pt, dashed] (4,-0.65) -- (7,-0.65);
  \draw[BookTeal, line width=0.9pt] (7,-1.0) rectangle (12,-0.3);
  \node at (9.5,-0.65) {B 驱动};
  \node[hw label, anchor=east] at (-0.2,-0.65) {总线};
  \node[hw label] at (5.5,-1.2) {高阻};
  \draw[BookMuted, line width=0.6pt, <->] (4,2.42) -- node[hw label, above]{先断后通的交断间隙} (7,2.42);
\end{pdfdiagram}
\diagramnote{上：两个三态驱动者同时使能时的冲突通路——电流从 VCC 经 A 的上拉管、总线、B 的下拉管直到 GND（图中只画出导通的管子，另一只方向的管子均关断），按每个输出级导通电阻约 25$\Omega$ 估算约 100mA，总线电平停在中点附近，既不是合法高也不是合法低。下：正确的驱动权交接——A 的使能先撤、输出进入高阻，留出先断后通的切换间隙后 B 才使能，任何时刻总线上最多一个驱动者。}

争用是一个方向的故障，另一个方向的故障是无人驱动。若某一时刻所有使能都无效，总线悬空，电压会被漏电和耦合慢慢推走——这正是本章开篇悬空按钮问题的总线版本。工程上常用总线保持器兜底：它是挂在总线上的一个弱反馈锁存，最后一个驱动者留下什么电平，它就以微弱电流维持什么电平；新驱动者接管时，其强驱动轻易盖过这个弱保持，总线正常翻转。注意分工：总线保持器解决“无人驱动时漂移”，完全不解决“多人驱动时争用”——使能互斥仍是设计必须证明的性质，不能交给器件去兜底。

\begin{longtable}{L{0.24\linewidth}L{0.66\linewidth}}
\toprule
关键问题 & 预期结果 \\
\midrule
使能是否互斥 & 任何时刻至多一个 \code{OE_n} 有效，由译码器 one-hot 或仲裁逻辑保证 \\
驱动者切换间隙 & 旧驱动者先进入高阻，新驱动者再使能，留出先断后通的切换间隙 \\
无人驱动时 & 总线保持器或上拉/下拉电阻给出确定电平，不悬空 \\
上电与复位期间 & 所有 \code{OE_n} 默认无效，总线默认空闲，不误导后级 \\
\bottomrule
\end{longtable}

后文讲整机时会看到，存储器、外设和 CPU 的数据引脚几乎都是三态的；一张原理图上几十根数据线的和平共处，靠的就是这里写下的互斥约定。

\topic{开漏输出与线与}
三态输出的思路是“要么全力驱动，要么彻底退出”。开漏输出走另一条路：它只保留下拉管，干脆没有上拉管。输出管导通时，线被强拉到 0；输出管关断时，驱动器松手——注意不是输出 1，而是什么都不做。高电平由线路上唯一的公共上拉电阻提供：所有驱动器都松手后，电流经电阻向线路电容充电，电压回升到电源。

这个结构天然允许一条线挂多个驱动者。任何一个驱动者拉低，全线就是 0；只有全部松手，线才回到 1。把“松手”读成赞成、“拉低”读成否决，线上的电平就是全体的 AND——没有门的实体，连线本身完成了与运算，所以叫线与。与三态方案对比，差别很清楚：三态要求使能互斥，违反就是电源到地的争用事故；开漏不存在这个事故，最坏情况只是多个驱动者同时拉低，线依然可靠为 0。代价集中在高电平一侧：上升沿靠电阻慢慢充电，快不起来；拉低期间电阻上持续流过电流，白白耗电。

I2C 总线的 SDA 和 SCL 就是这个模型。每个设备都只能拉低或松手，起始、停止、应答全部是“谁在何时拉低”的约定；从机处理不过来时按住 SCL 不放，主机发现时钟线升不回去，就知道必须等待——这就是时钟延展。这类总线的动作电平因此都是低有效，与前文低有效信号约定出自同一个电气原因：低电平是有源强驱动，可靠且谁都能发起；高电平是电阻慢慢充出来的，只配当默认态。

上拉电阻的取值是开漏设计的主要权衡，它同时卡在两个约束之间。上升沿一侧：时间常数 $\tau = R \cdot C_{\mathrm{bus}}$，电阻越大，回到高电平越慢。拉低一侧：电流 $I = V_{\mathrm{DD}} / R$，电阻越小，拉低时耗电越多，且驱动器的吸收电流能力有限，电流超限会把低电平抬高、侵蚀低电平余量。以 $V_{\mathrm{DD}} = 3.3\mathrm{V}$、总线电容 100pF 为例：

\begin{longtable}{L{0.14\linewidth}L{0.22\linewidth}L{0.24\linewidth}L{0.3\linewidth}}
\toprule
上拉电阻 & 拉低电流 $V_{\mathrm{DD}}/R$ & 时间常数 $\tau = R \cdot C$ & 适用场合 \\
\midrule
10 k$\Omega$ & 0.33 mA & 1000 ns & 低速、省电优先的场合 \\
4.7 k$\Omega$ & 0.70 mA & 470 ns & I2C 标准模式（100 kHz）常见取值 \\
2.2 k$\Omega$ & 1.5 mA & 220 ns & I2C 快速模式（400 kHz）常见取值 \\
1.0 k$\Omega$ & 3.3 mA & 100 ns & 边沿更快，但已略超常见的 3 mA 拉低限值 \\
\bottomrule
\end{longtable}

电压升到约 70\% 需要 $1.2\tau$：4.7 k$\Omega$ 时约 560 ns，尚在标准模式对上升沿的容限之内；总线电容更大或速率更高，就必须减小电阻，或者改用有源上拉电路。读表的方向不要反：不是“电阻越小越好”，而是在上升沿速度和拉低电流之间取交点。

到这里，三种输出级形态可以放在一起读：推挽输出任何时候都独占线路；三态输出靠互斥使能分时共享；开漏输出靠“只能拉低”让多驱动者和平共存。它们回答的是同一个问题——谁驱动这条线、何时驱动、没人驱动时谁来兜底。后文译码器产生的互斥 one-hot 信号，正是三态总线最常见的值班表；而开漏线把“兜底”直接做进了电气结构里。

译码器则常用于“编号变成单独使能”。例如 2-bit 编号选择 4 个寄存器中的一个写入，译码器输出 one-hot 写使能：同一时刻只允许一个目标被写。若译码器错误地产生两个 1，就可能同时写两个寄存器；若所有输出都是 0，本周期写操作会消失。

\begin{longtable}{L{0.18\linewidth}L{0.36\linewidth}L{0.36\linewidth}}
\toprule
结构 & 典型用途 & 必须保证 \\
\midrule
选择器 & 多个来源选一个输出 & 同一时刻只有一路被选中成为结果 \\
译码器 & 编号展开成 one-hot 使能 & 合法编号只激活一个目标 \\
编码器 & 多个请求压缩成编号 & 多请求同时有效时要有优先级或报错 \\
比较器 & 判断相等、大小或条件成立 & 输入解释方式必须明确 \\
\bottomrule
\end{longtable}

编码器和比较器也是常见组合电路。编码器把 one-hot 或优先级请求压缩成编号；比较器判断两个值是否相等，或者按某种解释比较大小。安全启动控制板可以用简单编码器把故障来源压缩成故障码：没有故障为 00，温度故障为 01，电流故障为 10，若两者同时故障则输出优先级最高的一项或输出“多故障”码。关键在于规则必须写清楚，不能让两个输入同时为 1 时输出处于未定义状态。

\begin{longtable}{L{0.18\linewidth}L{0.18\linewidth}L{0.24\linewidth}L{0.28\linewidth}}
\toprule
\code{temp_alarm} & \code{current_alarm} & 故障码 & 说明 \\
\midrule
0 & 0 & 00 & 无故障 \\
1 & 0 & 01 & 温度故障 \\
0 & 1 & 10 & 电流故障 \\
1 & 1 & 11 & 多故障或最高优先级故障 \\
\bottomrule
\end{longtable}

若系统选择优先级编码，就要把优先级写进规格。例如温度故障比电流故障优先，则两个报警同时为 1 时输出温度故障码；若系统更重视诊断完整性，则两个报警同时为 1 时输出多故障码。二者没有绝对优劣，但不能在电路里含糊。

\begin{longtable}{L{0.2\linewidth}L{0.34\linewidth}L{0.36\linewidth}}
\toprule
编码策略 & 规则 & 适合场景 \\
\midrule
优先级编码 & 多个请求同时有效时选择最高优先级 & 中断控制、快速响应最重要的请求 \\
多故障编码 & 多个请求同时有效时输出专用多故障码 & 故障诊断和维护记录 \\
非法状态报警 & 非 one-hot 输入触发错误输出 & 希望尽早暴露编码器前级错误 \\
\bottomrule
\end{longtable}

\topic{译码器、编码器与优先编码器}
前文把译码器和编码器各用一段话带过，这两种结构值得展开，因为它们是“编号”与“选择”之间往返转换的标准件。

译码方向先一般化：n 位输入、$2^n$ 条输出，第 k 条输出正好是最小项 $m_k$——它只在输入等于 k 时为 1。前文的 2 到 4 译码器是 n=2 的特例；常见的 3 到 8 译码器（74HC138 一类）还带使能端，使能无效时所有输出一律无效。使能端给了译码器两个能力：级联，以及给片选加一个总开关。译码器的关键性质是输出互斥：合法输入下恰好一条输出有效。这个 one-hot 性质正好满足三态总线和寄存器写使能要求的“唯一驱动者”——译码器因此是“地址变片选”的通用件，存储器选体、外设选择、寄存器堆写入都靠它。

编码器走反方向：$2^n$ 条 one-hot 输入，输出 n 位编号，即“第几条线有效”。普通编码器要求输入严格 one-hot；若两条输入同时有效，输出没有定义。这不是电路的缺陷，而是输入约定被打破——编码器只承诺在约定之内给出正确编号。

优先编码器把这条约定放宽：允许多个输入同时有效，优先级最高者胜出，并额外给出一个标志 G，表示“当前至少有一个请求有效”。以 4 到 2 优先编码器为例，约定 I3 优先级最高、I0 最低：

\begin{longtable}{L{0.26\linewidth}L{0.22\linewidth}L{0.42\linewidth}}
\toprule
I3 I2 I1 I0 & Y1 Y0 \quad G & 说明 \\
\midrule
0\quad 0\quad 0\quad 0 & 0\quad 0\quad 0 & 无请求，G=0 \\
0\quad 0\quad 0\quad 1 & 0\quad 0\quad 1 & 只有 I0 请求 \\
0\quad 0\quad 1\quad X & 0\quad 1\quad 1 & I1 在场，I0 被压住 \\
0\quad 1\quad X\quad X & 1\quad 0\quad 1 & I2 在场，低两位被压住 \\
1\quad X\quad X\quad X & 1\quad 1\quad 1 & I3 一票定音 \\
\bottomrule
\end{longtable}

表中 X 表示该输入不影响本行输出。逐列读出布尔式：

\begin{lstlisting}
Y1 = I3 OR I2
Y0 = I3 OR (NOT I2 AND I1)
G  = I3 OR I2 OR I1 OR I0
\end{lstlisting}

验算两行：I2=1、I3=0 时，Y1=1、Y0=0，编号 10，且无论 I1、I0 如何都不变——这正是“压住”的含义；只有 I1=1 时，Y1=0、Y0=1，编号 01。G 与编号相互独立：全 0 输入时编号也是 00，全靠 G 区分“有请求、编号为 00”和“根本没有请求”。

优先编码器是中断请求排队的雏形。多个设备共用一路请求通道时，同时请求是常态而不是异常；固定优先级的编码器先报出最高优先级设备的编号，G 则直接充当“有待处理请求”的总信号。后文中断控制器会在此基础上增加屏蔽、嵌套和可编程优先级，但裁决的第一步就是这里这几行式子。

两种结构的级联也值得对照。译码器靠使能端扩展：两个 3 到 8 译码器，地址最高位 A3 直接接高位片的使能、反相后接低位片的使能，就得到 4 到 16 译码，输出仍然互斥。优先编码器靠 G 串联扩展：高位片的 G 决定低位片的编号是否被采用，全局优先级仍然唯一。选择器的树形扩展、译码器的使能级联、优先编码器的 G 串联，是同一个思想的三种形态：用少量控制端，把“同一时刻只有一个胜者”的性质扩展到大网络。

比较器也必须说明数值解释。同样的 bit 模式，用无符号数解释和用补码有符号数解释，大小关系可能不同。以 4-bit 为例，\code{1111} 作为无符号数是 15；作为补码有符号数是 -1。若比较器输出 \code{less_than}，规格必须说明它比较的是无符号大小、有符号大小，还是只比较 bit 模式是否相等。

\begin{longtable}{L{0.2\linewidth}L{0.28\linewidth}L{0.42\linewidth}}
\toprule
比较类型 & 示例 & 说明 \\
\midrule
相等比较 & \code{A == B} & 只关心每一位是否相同，通常不涉及数值解释 \\
无符号大小 & \code{1111 > 0001} & \code{1111} 表示 15 \\
有符号大小 & \code{1111 < 0001} & \code{1111} 表示 -1，\code{0001} 表示 1 \\
条件比较 & 报警级别是否超过阈值 & 必须说明编码范围和非法值处理 \\
\bottomrule
\end{longtable}

\topic{数值比较器与四变量卡诺图案例}
相等比较只问“每位是否相同”，数值比较器还要多问一句“谁大”。裁决规则与字典序相同：从最高位开始逐位比较，第一对不同的位决定全局结果，更低的位不再有机会发言。以 4 位无符号数为例，先逐位求相等信号 \code{eq_i = A_i XNOR B_i}，再让高位的一票定音权逐级传递：

\begin{lstlisting}
gt = (A3 AND NOT B3)
  OR (eq3 AND A2 AND NOT B2)
  OR (eq3 AND eq2 AND A1 AND NOT B1)
  OR (eq3 AND eq2 AND eq1 AND A0 AND NOT B0)

eq = eq3 AND eq2 AND eq1 AND eq0
lt = 与 gt 对称，A、B 互换即可
\end{lstlisting}

读 \code{gt} 的每一项：第一项说“最高位就分出胜负”；第二项说“最高位打平，则由次高位裁决”，依此类推。\code{eq_i} 在这里身兼两职——既是相等比较的输出，又是大小比较中“把裁决权交给下一位”的通行证。这与逐位进位链在结构上同构：都是高位优先、逐级传递，级数同样随位宽线性增长。

单片 4 位比较器（74HC85 一类）还带三个级联输入：\code{I(A>B)}、\code{I(A=B)}、\code{I(A<B)}。它们的作用是：当本片四位全部打平时，比较结果不由本片决定，而是透传级联输入。把低位片的三个输出接到高位片的级联输入，两片就串成 8 位比较器，四片串成 16 位——裁决权从最低位一路传递到最高位，与式子里的 \code{eq} 链严格对应。有符号与无符号的差别集中在最高有效位的解释上（补码的最高位是负权），级联结构本身不变。

比较器用到的“相等信号加传递链”是一种标准手法。下面换一个完整演算，把前文只在文字层面讲过的卡诺图，在一个含无关项的四变量函数上走完全程。题目：检测 8421 BCD 码是否大于等于 5。输入 A、B、C、D（A 为最高位），按约定只会出现 0 到 9；10 到 15 是无关项，记为 d。

先列完整真值表，一行都不跳：

\begin{longtable}{L{0.16\linewidth}L{0.24\linewidth}L{0.12\linewidth}L{0.38\linewidth}}
\toprule
十进制 & A B C D & F & 说明 \\
\midrule
0 & 0 0 0 0 & 0 & 小于 5 \\
1 & 0 0 0 1 & 0 & 小于 5 \\
2 & 0 0 1 0 & 0 & 小于 5 \\
3 & 0 0 1 1 & 0 & 小于 5 \\
4 & 0 1 0 0 & 0 & 小于 5 \\
5 & 0 1 0 1 & 1 & 大于等于 5 \\
6 & 0 1 1 0 & 1 & 大于等于 5 \\
7 & 0 1 1 1 & 1 & 大于等于 5 \\
8 & 1 0 0 0 & 1 & 大于等于 5 \\
9 & 1 0 0 1 & 1 & 大于等于 5 \\
10 & 1 0 1 0 & d & 无关项：约定不会出现 \\
11 & 1 0 1 1 & d & 无关项 \\
12 & 1 1 0 0 & d & 无关项 \\
13 & 1 1 0 1 & d & 无关项 \\
14 & 1 1 1 0 & d & 无关项 \\
15 & 1 1 1 1 & d & 无关项 \\
\bottomrule
\end{longtable}

再填四变量卡诺图。A、B 沿行方向按 Gray 码排列，C、D 沿列方向按 Gray 码排列，每个格子填入对应最小项的函数值：

\begin{longtable}{L{0.22\linewidth}L{0.14\linewidth}L{0.14\linewidth}L{0.14\linewidth}L{0.14\linewidth}}
\toprule
AB \textbackslash CD & 00 & 01 & 11 & 10 \\
\midrule
00 & 0 & 0 & 0 & 0 \\
01 & 0 & 1 & 1 & 1 \\
11 & d & d & d & d \\
10 & 1 & 1 & d & d \\
\bottomrule
\end{longtable}

圈组的规则只有两条：每个圈必须是 $2^k$ 个相邻格子组成的矩形，且每个 1 至少被一个圈覆盖；无关项可以当 1 圈进去，也可以当 0 留在圈外，取舍的唯一标准是能否把圈撑得更大。按这个标准可以圈出三组。第一组：A=1 的下面两行共 8 格（m8 到 m15），六个无关项全部圈入，得到单文字项 A——B、C、D 在圈内部取遍两种值，全部消去。第二组：B=1 且 C=1 的四格（m6、m7、m14、m15），其中两格是无关项，得到 BC。第三组：B=1 且 D=1 的四格（m5、m7、m13、m15），同样借两个无关项撑大，得到 BD。三个圈合起来：

\begin{lstlisting}
F = A OR (B AND C) OR (B AND D) = A OR (B AND (C OR D))
\end{lstlisting}

逐行验算：5=0101 由 BD 覆盖，6=0110 由 BC 覆盖，7 同时落在两个圈（BC 和 BD）里（重复覆盖不改变函数），8、9 由 A 覆盖；0 到 4 各行 A=0，且 B、C、D 不满足任何一项，输出 0。全部符合真值表。

为什么这样选圈，而不是逐对合并？对照不用无关项的化简就清楚了。若把 d 全部当 0，m8、m9 只能两两合并成 \code{A AND (NOT B) AND (NOT C)}，m5、m7 合并成 \code{(NOT A) AND B AND D}，m6、m7 合并成 \code{(NOT A) AND B AND C}——同样是三个乘积项，但每项有三个文字，合计 9 个文字；借入无关项后，三项合计只有 5 个文字，门更小、扇入更低。无关项的价值就在这里：它不增加任何需要覆盖的 1，却允许把已有的 1 圈进更大的组里。

最后必须重申前文对无关项的警告。这个化简把 10 到 15 实现成了 1——在“输入只会是合法 BCD 码”的约定下无害；一旦输入可能来自外部连接器或故障状态，这个约定就可能被打破，届时 d 必须改回确定值重新化简。化简服从系统约定，而不是相反。

现在把贯穿案例写成一个完整组合网络。安全启动控制板至少需要两个输出：\code{fault_lamp} 和 \code{allow_start}。故障灯由报警信号决定；启动允许由启动按钮、安全条件和无故障共同决定。

\begin{lstlisting}
fault_lamp = temp_alarm OR current_alarm
no_fault = NOT fault_lamp
allow_start = start_request AND cover_closed
              AND emergency_ok AND no_fault
\end{lstlisting}

这组表达式可以继续展开成门网络。第一层 OR 汇总故障信号；第二层 NOT 得到无故障条件；第三层多输入 AND 汇总启动按钮和安全条件。若硬件库里没有四输入 AND，可以用二输入 AND 逐级合成：

\begin{lstlisting}
safe_chain = cover_closed AND emergency_ok
request_ok = start_request AND no_fault
allow_start = safe_chain AND request_ok
\end{lstlisting}

这种重写在逻辑上等价，但硬件上改变了门级分组。若 \code{cover_closed} 和 \code{emergency_ok} 来自同一个连接器，先合成 \code{safe_chain} 可能便于布线；若 \code{fault_lamp} 驱动多个后级，\code{no_fault} 前后可能需要缓冲。组合逻辑的表达式只是第一层结果，门级分组、负载和时序仍要继续检查。

若器件库只允许二输入 NAND，还可以把同一逻辑展开为 NAND-only 网络。每个 AND 由“一次 NAND 得到反相信号，再用 NAND 自反相恢复”为正逻辑。中间的反相信号必须命名清楚，否则很容易在后续维护中把 \code{safe_chain_n} 当作 \code{safe_chain} 使用。

\begin{lstlisting}
n1 = cover_closed NAND emergency_ok
safe_chain = n1 NAND n1

n2 = start_request NAND no_fault
request_ok = n2 NAND n2

n3 = safe_chain NAND request_ok
allow_start = n3 NAND n3
\end{lstlisting}

NAND-only 写法证明了功能可实现，但不自动证明它就是最好的实现。线性串接、平衡二输入树和 NAND-only 分解在稳定真值表上可以等价，门级层数、局部负载、中间节点数量和毛刺形态却可能不同。平衡树通常缩短最长逻辑层数，线性结构可能更便于让某个晚到信号放在最后一级，NAND-only 结构可能更符合某些标准门芯片资源。正式规格应保留可审查的安全表达式，再单独说明实际门级映射，避免把实现技巧误写成需求本身。

顺这条链再检查一遍内部信号：\code{fault_lamp} 表示任一报警成立，OR 路径要能驱动指示灯和后级逻辑；\code{no_fault} 是它的反相，必须由反相器恢复为强 0/1；\code{safe_chain} 要求保护盖和急停都满足，任一断开时默认禁止；\code{request_ok} 要求有启动请求且无故障，按钮抖动不能传到后级；\code{allow_start} 是全部条件的汇合，后级何时采样、有无毛刺风险都要写明；\code{motor_enable_cmd} 是发给执行机构的命令，必须经过时序边界或专用安全链路，不能直接把组合输出接出去。

组合逻辑设计到这里仍然没有引入“步骤执行”。所有输入同时作用于网络，所有中间信号同时传播。若 \code{temp_alarm} 从 0 变 1，\code{fault_lamp} 会在 OR 路径延迟后变为 1，\code{no_fault} 会在反相器延迟后变为 0，\code{allow_start} 会在 AND 路径延迟后关闭。这个传播链条不是程序调用栈，而是多条电气路径的延迟叠加。

\code{motor_enable_cmd} 需要比 \code{allow_start} 更谨慎。前者是给外部执行机构的命令，后者只是内部组合判断。直接令 \code{motor_enable_cmd = allow_start} 在最小示例里看似合理，但正式系统通常会增加时序边界、复位保持、驱动级和独立停机链路。原因很简单：组合判断可以短暂过渡，执行机构不应响应这种过渡。

\begin{lstlisting}
allow_start = start_request AND cover_closed
              AND emergency_ok AND no_fault

motor_enable_cmd = registered_allow_start AND power_stage_ready
\end{lstlisting}

上式中的 \code{registered_allow_start} 表示已经在时钟边界保存过的允许结果，\code{power_stage_ready} 表示功率级自身已准备好并未报告故障。数字硬件基础部分还没有正式引入触发器，所以这里只把它作为边界预告：组合逻辑负责给出当前条件下的判断，执行机构命令通常还需要由时序逻辑和驱动电路接管。这样组织，能把“判断是否允许”和“真正驱动负载”分成两个可独立检查的层次。

可以把安全启动控制板整理成一份组合规格。规格不必一开始就写门级细节，但必须把输入语义、输出语义和异常策略写清楚：输入前提是板外触点已经完成默认值、电平和消抖处理，原始触点不直接进入安全组合式；允许条件是启动请求、保护盖闭合、急停正常、无故障四者缺一不可；故障输出要求温度或电流报警都点亮故障灯，同时报警要有明确编码策略；默认安全要求未知、断线、非法编码一律倒向禁止启动，化简逻辑不得破坏这条策略；使用时刻要求后级只在信号稳定后使用，异步使能必须额外做毛刺检查。这样做的好处是，之后无论用分立门、可编程逻辑，还是把同一逻辑放进更大的控制器，行为边界都保持一致。

下面给出这类小组合模块的完整书写顺序。它看起来比直接写一个公式更长，但每一步都对应一种常见错误：信号没定义、非法状态没处理、公式和需求不一致、实现后电气条件不满足。第一步，列出原始输入、内部语义信号和输出，不把物理触点直接当作可靠 bit。第二步，说明每个信号的有效电平和默认值，避免低有效、断线和上电默认错误。第三步，写出正常行为和异常行为，不让非法组合默默进入允许状态。第四步，写逻辑表达式或真值表，让需求可以机械检查。第五步，映射到门、选择器、译码器或编码器，明确驱动来源和选择关系。第六步，检查延迟、负载、毛刺、电平余量，证明实物能按时产生可靠输出。第七步，记录测量或仿真结果，让设计从文字规格变成可验证的对象。

把这七步用于 \code{allow_start}，可以得到更严格的文字规格：输入必须已经完成调理；\code{start_request} 为 1 表示一次经过消抖或同步确认的启动请求，缺失时不启动；\code{cover_closed} 为 1 表示保护盖闭合且传感器链路有效，缺失时不启动；\code{emergency_ok} 为 1 表示急停链路处于允许运行状态，缺失时不启动或停机；\code{no_fault} 为 1 表示没有温度或电流报警，缺失时不启动并报警；组合路径已越过最坏延迟，输出才算稳定，此前后级暂不采样、也不驱动执行机构。

若某个条件缺失却仍能启动，就是需求错误或实现错误；若输出尚未稳定就驱动执行机构，就是时序边界错误。组合逻辑的严肃性正在于此：它没有软件中的异常处理栈，只有电路在每一种输入情况下实际产生的电平。

为了让贯穿案例形成闭环，可以把安全启动控制板从外部物理输入一路写到执行命令。该表不是新增功能，而是把前面分散的原则合成一条可审查链路：原始线缆先进入电气边界，调理后得到内部语义信号，组合网络只处理已经可靠的 bit，危险输出再经过时序和驱动边界。

\begin{longtable}{L{0.17\linewidth}L{0.28\linewidth}L{0.45\linewidth}}
\toprule
层次 & 典型对象 & 约定要求 \\
\midrule
原始输入 & \code{start_button_raw}、传感器触点、报警线 & 允许抖动和噪声，但必须有默认路径、保护和极性说明 \\
输入调理 & 上拉/下拉、滤波、施密特触发、电平转换 & 把板外现象整理成稳定内部电平，不让悬空和慢边沿直接进入核心 \\
语义信号 & \code{start_request}、\code{cover_closed}、\code{emergency_ok} & 每个 0/1 含义清楚，低有效信号先转换为可审查语义 \\
组合判断 & \code{fault_lamp}、\code{no_fault}、\code{allow_start} & 真值表覆盖正常、异常、非法和默认状态，化简不破坏安全策略 \\
时序边界 & \code{registered_allow_start}、复位保持、采样时刻 & 只在组合路径稳定后保存或使用结果，避免毛刺进入危险动作 \\
输出驱动 & \code{motor_enable_cmd}、指示灯驱动、功率级使能 & 具备足够驱动能力，上电和复位期间保持禁止，故障时进入安全状态 \\
验证记录 & 原理图、数据手册、示波器、逻辑分析仪、测试矩阵 & 说明电平、延迟、负载、毛刺和异常输入处理都满足约定 \\
\bottomrule
\end{longtable}

这张闭环表也说明为什么数字硬件基础部分要同时讲电平、器件、门和组合逻辑。如果只看组合表达式，会把 \code{start_request} 当作天然可靠的变量；如果只看器件，又难以说明为何这些输入最终必须合成为一个 \code{allow_start}。正式硬件说明应当让两者互相校验：每个语义信号都能追溯到电气边界，每个组合输出都能追溯到真值表和电流路径，每个执行命令都能追溯到时序和驱动电路。

\topic{端到端示例：用两片 74HC00 实现安全启动控制板}
前面所有片段——真值表、NAND-only 展开、芯片引脚、噪声余量——现在可以合成一次完整演练。目标是让 \code{allow_start} 从自然语言需求一路走到可焊接的引脚分配。

第一步，需求表达式（安全语义层，保持可审查）：

\begin{lstlisting}
fault_lamp = temp_alarm OR current_alarm
no_fault = NOT fault_lamp
allow_start = start_request AND cover_closed
              AND emergency_ok AND no_fault
\end{lstlisting}

第二步，NAND-only 展开（实现层，每个反相中间信号命名清楚）。\code{fault_lamp} 与 \code{no_fault} 用一路 OR 芯片（如 74HC32）或二极管 OR 加缓冲实现即可，这里聚焦 AND 链：

\begin{lstlisting}
n1 = cover_closed NAND emergency_ok
safe_chain = n1 NAND n1
n2 = start_request NAND no_fault
request_ok = n2 NAND n2
n3 = safe_chain NAND request_ok
allow_start = n3 NAND n3
\end{lstlisting}

第三步，数门。共 6 个二输入 NAND；一片 74HC00 有 4 路，需要两片。74HC00 的封装是 14 脚双列直插：左边 1–7 脚，右边 14–8 脚，缺口朝上时左上角是 1 脚。

\begin{circuit}[x=1cm,y=1cm]
  \draw[line width=0.6pt, rounded corners=2pt] (0,0) rectangle (4.2,4.55);
  \draw[line width=0.6pt] (1.9,4.55) arc (180:360:0.2);
  \node[hw label] at (2.1,2.3) {74HC00（U1）};
  \foreach \y/\p/\f/\d in {3.9/1/{1A}/{$\leftarrow$ cover\_closed}, 3.25/2/{1B}/{$\leftarrow$ emergency\_ok}, 2.6/3/{1Y}/{$\rightarrow$ n1}, 1.95/4/{2A}/{$\leftarrow$ n1}, 1.3/5/{2B}/{$\leftarrow$ n1}, 0.65/6/{2Y}/{$\rightarrow$ safe\_chain}, 0.0/7/{GND}/{GND}} {
    \draw[line width=0.5pt] (-0.4,\y+0.32) -- (0,\y+0.32);
    \node[hw label, anchor=east] at (-0.55,\y+0.32) {\p\ \d};
    \node[hw label, anchor=west] at (0.15,\y+0.32) {\f};
  }
  \foreach \y/\p/\f/\d in {3.9/14/{VCC}/{VCC}, 3.25/13/{4B}/{$\leftarrow$ n2}, 2.6/12/{4A}/{$\leftarrow$ n2}, 1.95/11/{4Y}/{$\rightarrow$ request\_ok}, 1.3/10/{3B}/{$\leftarrow$ start\_request}, 0.65/9/{3A}/{$\leftarrow$ no\_fault}, 0.0/8/{3Y}/{$\rightarrow$ n2}} {
    \draw[line width=0.5pt] (4.2,\y+0.32) -- (4.6,\y+0.32);
    \node[hw label, anchor=west] at (4.75,\y+0.32) {\p\ \d};
    \node[hw label, anchor=east] at (4.05,\y+0.32) {\f};
  }
\end{circuit}
\diagramnote{第一片 74HC00 的引脚分配：门 1 产生 \code{n1}，门 2 把 \code{n1} 自反相得到 \code{safe_chain}，门 3 产生 \code{n2}，门 4 把 \code{n2} 自反相得到 \code{request_ok}。缺口（上方半圆）决定引脚编号方向——原理图和实物靠它对应。}

第二片（U2）完成最后两级：门 1 计算 \code{n3 = safe_chain NAND request_ok}，门 2 自反相得到 \code{allow_start}；门 3、门 4 未使用。未使用的 CMOS 输入必须接确定电平（例如都接 VCC），否则悬空输入会让输出级随机翻转、白白耗电。

\begin{longtable}{L{0.16\linewidth}L{0.44\linewidth}L{0.3\linewidth}}
\toprule
资源 & 分配 & 检查点 \\
\midrule
U1 门 1–2 & \code{cover_closed}、\code{emergency_ok} 合成 \code{safe_chain} & 中间信号 \code{n1} 是反相电平，不可误用 \\
U1 门 3–4 & \code{start_request}、\code{no_fault} 合成 \code{request_ok} & \code{no_fault} 来自故障汇总级 \\
U2 门 1–2 & \code{safe_chain}、\code{request_ok} 合成 \code{allow_start} & 输出送时序边界，不直接驱动执行机构 \\
U2 门 3–4 & 未使用 & 输入接确定电平，输出悬空不理 \\
电源 & 14 脚 VCC、7 脚 GND，两片各自就近 0.1$\mu$F 去耦 & 去耦电容贴近电源脚 \\
\bottomrule
\end{longtable}

第四步，电气参数核算。以 74HC00 数据手册在 4.5V 供电下的保证值为例：输出高电平最差约 $V_\mathrm{OH}=3.84$V（4mA 负载），输出低电平最差 $V_\mathrm{OL}=0.33$V；输入高阈值 $V_\mathrm{IH}=3.15$V，输入低阈值 $V_\mathrm{IL}=1.35$V。逐级核算：

\begin{longtable}{L{0.3\linewidth}L{0.28\linewidth}L{0.32\linewidth}}
\toprule
检查项 & 算式 & 结论 \\
\midrule
高电平余量 & $3.84\mathrm{V} - 3.15\mathrm{V}$ & 约 0.69V，合格 \\
低电平余量 & $1.35\mathrm{V} - 0.33\mathrm{V}$ & 约 1.0V，合格 \\
级联三级后 & 每级输出都重新驱动到 3.84V/0.33V & CMOS 恢复电平，余量不逐段衰减 \\
\bottomrule
\end{longtable}

注意第三行与二极管 OR 的本质差别：二极管每级吃掉 0.7V，三级后余量所剩无几；CMOS 门每级都把电平重新驱动到接近电源和地，余量被\hwkey{恢复}——这就是为什么通用逻辑网络必须用有源门，而不是二极管逻辑。

至此，\code{allow_start} 可以回答本章开头的全部追问：需求来自真值表，实现落在两片具体芯片的 12 个引脚上，电气有余量核算，时序上它只送到时序边界而不是执行机构。这个从自然语言到引脚的完整链条，就是本章要达到的目标。

\topic{端到端示例二：三级报警优先权电路}
上一例回答的是“所有安全条件同时成立才允许动作”，还有另一类同样常见的需求：多个请求可能同时到来，必须按优先级裁决，高优先级抢占输出并屏蔽低优先级。本例用一片 74HC04、一片 74HC00 和一片 74HC32，把一台设备的三级报警面板从需求走到芯片分配，流程与上一例完全相同。

需求如下。面板有三个请求源：火警 \code{fire}（最高优先级）、故障 \code{fault}、提示 \code{info}（最低优先级），对应三盏指示灯 \code{fire_lamp}、\code{fault_lamp}、\code{info_lamp}。任一时刻只允许点亮当前最高优先级的那一盏：高优先级请求出现时立即抢占，并把低优先级输出强制为 0。注意屏蔽不是“盖住显示”，而是让低优先级输出在逻辑上就是 0——这样无论后级接的是灯、蜂鸣器还是记录电路，看到的都是被裁决过的结果。全部请求消失时三盏都熄灭。

先列完整真值表。三个请求共八种组合，裁决结果如下：

\begin{longtable}{L{0.28\linewidth}L{0.32\linewidth}L{0.3\linewidth}}
\toprule
请求（fire, fault, info） & 输出（fire/fault/info\_lamp） & 裁决结果 \\
\midrule
0, 0, 0 & 0, 0, 0 & 无请求，全部熄灭 \\
0, 0, 1 & 0, 0, 1 & 只有提示 \\
0, 1, 0 & 0, 1, 0 & 故障点亮，提示未请求 \\
0, 1, 1 & 0, 1, 0 & 故障屏蔽提示 \\
1, 0, 0 & 1, 0, 0 & 火警独占 \\
1, 0, 1 & 1, 0, 0 & 火警屏蔽提示 \\
1, 1, 0 & 1, 0, 0 & 火警屏蔽故障 \\
1, 1, 1 & 1, 0, 0 & 火警屏蔽全部 \\
\bottomrule
\end{longtable}

从真值表直接读出化简结果。\code{fire_lamp} 在 \code{fire}=1 的四行全为 1，与另两个输入无关，所以 \code{fire_lamp = fire}——最高优先级输出不经过任何门，路径最短，这正符合它的安全地位。\code{fault_lamp} 只在 (0,1,0) 与 (0,1,1) 两行为 1，化简得 \code{fault AND (NOT fire)}：火警的反相信号作为一个输入项，它一来，故障灯就被逻辑本身关断。\code{info_lamp} 只在 (0,0,1) 一行为 1，即 \code{info AND (NOT fire) AND (NOT fault)}；用德摩根律改写成 \code{info AND NOT(fire OR fault)}，可以省下一级三输入 AND：

\begin{lstlisting}
fire_n     = NOT fire
m1         = fault NAND fire_n      -- 74HC00 门1
fault_lamp = NOT m1                 -- 74HC04
any_hi     = fire OR fault          -- 74HC32 门1
any_hi_n   = NOT any_hi             -- 74HC04
m2         = info NAND any_hi_n     -- 74HC00 门2
info_lamp  = NOT m2                 -- 74HC04
any_alarm  = any_hi OR info         -- 74HC32 门2，校验输出
fire_lamp  = fire                   -- 直连，不经门
\end{lstlisting}

芯片分配如下。74HC04（U1）用四路反相器，产生 \code{fire_n}、\code{fault_lamp}、\code{any_hi_n}、\code{info_lamp}；74HC00（U2）用两路 NAND，产生 \code{m1}、\code{m2}；74HC32（U3）用两路 OR，产生 \code{any_hi}、\code{any_alarm}。U1 门 5–6、U2 门 3–4、U3 门 3–4 未使用，输入一律接确定电平。其中 \code{any_alarm} 是故意多算的校验输出：它在逻辑上应恒等于 \code{fire OR fault OR info}，也应恒等于三盏灯输出的 OR——同一个事实沿两条独立路径各算一遍，上电测试时逐行比对，就是一种廉价的一致性检查。

\begin{longtable}{L{0.18\linewidth}L{0.44\linewidth}L{0.28\linewidth}}
\toprule
资源 & 分配 & 检查点 \\
\midrule
U1（74HC04）门 1–4 & \code{fire_n}、\code{fault_lamp}、\code{any_hi_n}、\code{info_lamp} & 门 5–6 未用，输入接确定电平 \\
U2（74HC00）门 1–2 & 产生 \code{m1}、\code{m2} & 门 3–4 未用，输入接确定电平 \\
U3（74HC32）门 1–2 & 产生 \code{any_hi}、\code{any_alarm} & 门 3–4 未用，输入接确定电平 \\
\code{fire} 直连 & \code{fire_lamp} 不经门 & 最高优先级路径最短；灯的电流由驱动级承担，不经过逻辑门 \\
裁决延迟 & \code{fire} 到 \code{fault_lamp} 三级门、到 \code{info_lamp} 四级门 & 延迟沿路径逐级累加，定量核算见第五节 \\
电源 & 每片 14 脚 VCC、7 脚 GND，就近 0.1$\mu$F 去耦 & 去耦电容贴近电源脚，理由见第五节专题 \\
一致性 & \code{any_alarm} 与三灯输出逐行比对 & 两条独立路径的结果必须一致 \\
\bottomrule
\end{longtable}

与上一例对照可以看到两点。第一，优先级不是额外的器件，而是写在真值表里的屏蔽关系：高优先级输入以反相形式出现在低优先级的乘积项里，它一来，低优先级输出就被逻辑本身关断。第二，优先级越高的输出经过的门越少——\code{fire_lamp} 零级、\code{fault_lamp} 最多三级、\code{info_lamp} 最多四级——“越重要的信号越快”在这里是自然结果，不需要刻意安排。至于三级门到底慢到什么程度，正是第五节要算的账。

把同一思路放大到 8086 和 80386，就能看出早期处理器为什么需要大量组合模块。以 Intel 公开产品资料所列参数为例，8086 于 1978 年引入，约 29,000 个晶体管，采用约 3 微米制造技术；Intel386 DX 于 1985 年引入，约 275,000 个晶体管，制造技术从约 1.5 微米进入后续 1 微米级别。这些数字本身不是本章重点，重点在于它们对应的结构变化：晶体管数量增加以后，可以容纳更宽的数据通路、更复杂的译码、更丰富的寄存器和更强的总线控制；工艺尺寸缩小以后，门和连线的寄生电容下降，单位距离连接更短，许多路径的延迟随之降低。

\begin{longtable}{L{0.16\linewidth}L{0.2\linewidth}L{0.24\linewidth}L{0.3\linewidth}}
\toprule
芯片 & 年代与规模 & 本章可观察的硬件变化 & 对延迟的意义 \\
\midrule
8086 & 1978 年，约 29,000 晶体管 & 16-bit 数据通路、指令译码、总线接口 & 已经需要大量组合控制与寄存器边界配合 \\
80386 & 1985 年，约 275,000 晶体管 & 32-bit 数据通路、更复杂保护和地址机制 & 更多片上逻辑减少外部协作等待，关键路径需重新组织 \\
现代 CPU & 数十亿级晶体管 & 多级缓存、乱序执行、预测、片上互连 & 不只让门更快，还减少取数和等待的次数 \\
\bottomrule
\end{longtable}

这里要避免一个误解：现代芯片低延迟不是因为“逻辑门不再有延迟”。门仍然有延迟，连线仍然有电容，存储器仍然需要访问时间。现代芯片做的是把延迟压低、隐藏、分层和局部化——这四个动词的具体含义放在本章第五节末尾展开，这里先记住结论：低延迟来自让常见情况少走远路，而不是让所有路径同样快。

\section{三、延迟、负载、毛刺与检查}

门电路在纸上像瞬间给出输出的函数，真实硬件却至少有四个限制。第一，输入变化后，输出需要传播延迟才稳定，组合逻辑层数越多，稳定越慢。第二，一个输出能驱动的后级输入数量有限，后级太多会让切换变慢，甚至让电平余量变差，大网络需要缓冲器和分层布线。第三，门的输出必须重新形成清楚的 0/1，不能只靠一个偏弱路径勉强改变节点电压。第四，多条路径延迟不同，输入变化过程中可能出现短暂毛刺，后级不能在毛刺期间采样。

考虑一个反例。若把十几个门的输入全接到同一个弱输出上，真值表仍然可以保持正确，但物理电路可能切换缓慢。输出从 0 变 1 时，要给所有后级输入电容充电；如果下一级在它尚未进入可靠高电平区间时读取，就可能看到旧值或不确定值。因此，分析组合逻辑时不能只数门的个数，还要考察最长路径、负载和稳定时间。

扇出问题可以放在安全启动控制板里理解。若 \code{fault_lamp} 只驱动一个小逻辑门，边沿可能很快；若它同时驱动指示灯缓冲、记录电路、禁止启动网络和外部连接器，等效负载明显增加。正确做法通常不是让一个弱节点直接驱动所有后级，而是分层缓冲：一个内部故障信号先驱动若干缓冲器，再由缓冲器分别驱动不同区域。缓冲器不改变逻辑含义，却改变驱动能力和负载隔离。同理，长线不能直接接到敏感输入上，容易耦合噪声和振铃，需要加强驱动、终端处理或同步采样；普通逻辑输出也不能直接承担指示灯负载，要用驱动级把指示灯的电流需求隔在逻辑电平之外。

短暂毛刺也很常见。假设一个输出由多条逻辑路径合成，输入同时变化时，每条路径稳定的时间可能不同。真值表只告诉我们变化前和变化后的稳定结果，却不保证中间不会短暂出现错误电平。如果后级只是另一个组合门，毛刺可能很快消失；如果后级在毛刺期间被时钟采样，错误就可能被保存下来。

一个典型毛刺表达式是：

\begin{lstlisting}
Y = (A AND B) OR ((NOT A) AND C)
\end{lstlisting}

假设 B=1、C=1。稳定状态下，不论 A=1 还是 A=0，Y 都应为 1：A=1 时第一项成立，A=0 时第二项成立。问题出现在 A 翻转的瞬间。若 A 从 1 变 0，第一条路径可能先关断，而 \code{NOT A} 经过反相器后才让第二条路径打开，中间可能出现一个短暂间隙，使 Y 从 1 短暂掉到 0。真值表看不见这个瞬间，因为真值表只描述稳定输入和稳定输出。

\begin{longtable}{L{0.18\linewidth}L{0.28\linewidth}L{0.44\linewidth}}
\toprule
时刻 & 路径状态 & Y 的风险 \\
\midrule
A=1 稳定 & \code{A AND B} 成立 & Y=1 \\
A 正在翻转 & 第一条路径可能已关，第二条路径尚未开 & Y 可能短暂为 0 \\
A=0 稳定 & \code{NOT A AND C} 成立 & Y=1 \\
\bottomrule
\end{longtable}

\begin{waveform}[x=1cm,y=1cm]
  \draw[hw bus] (0,0) -- (10.4,0);
  \node[hw label] at (10.4,-0.35) {时间};
  \node[hw label, anchor=east] at (-0.2,4.8) {A};
  \node[hw label, anchor=east] at (-0.2,3.6) {\code{NOT A}};
  \node[hw label, anchor=east] at (-0.2,2.0) {Y};
  \node[hw label, anchor=east] at (-0.2,0.4) {\code{Y\_safe}};
  \draw[BookInk, line width=1.0pt] (0,5.2) -- (3,5.2) -- (3,4.4) -- (10,4.4);
  \draw[BookMuted, line width=1.0pt] (0,3.2) -- (4,3.2) -- (4,4.0) -- (10,4.0);
  \draw[BookAccent, line width=1.1pt] (0,2.4) -- (3,2.4) -- (3,1.6) -- (4,1.6) -- (4,2.4) -- (10,2.4);
  \draw[BookTeal, line width=1.1pt] (0,0.8) -- (10,0.8);
  \draw[BookMuted, dashed, line width=0.4pt] (3,-0.3) -- (3,5.5);
  \draw[BookMuted, dashed, line width=0.4pt] (4,-0.3) -- (4,5.5);
  \draw[BookRed, line width=1.2pt] (3.15,1.35) -- (3.85,1.35);
  \node[hw label, text=BookRed] at (3.5,1.0) {毛刺};
  \node[hw label] at (3.5,6.2) {一个反相器延迟};
  \draw[hw bus, BookMuted] (3,5.6) -- (4,5.6);
\end{waveform}
\diagramnote{B=C=1 时，Y 的稳定值本应从 1 到 1 不变。但 A 翻转到 \code{NOT A} 生效之间隔着一个反相器延迟：上路先关、下路未开，Y 出现一个真值表看不见的窄 0 脉冲。加入共识项 \code{B AND C} 后（\code{Y\_safe}），有一条不依赖 A 的路径全程把输出钉在 1。}

若这个 Y 只是驱动另一个组合网络，并且后级只在更晚的时钟边界采样，毛刺可能不会造成状态错误；若 Y 直接作为写使能、复位、片选或异步控制信号，短暂毛刺就可能触发真实动作。安全启动控制板里的 \code{allow_start} 尤其需要谨慎：即使最终会回到正确值，也不能让一个窄脉冲误开电机使能。

上面的毛刺可以用增加共识项来消除。表达式 \code{Y = (A AND B) OR ((NOT A) AND C)} 在 B=1、C=1 且 A 翻转时有静态 1 毛刺风险。加入第三项 \code{B AND C} 后，A 翻转期间仍有一条不依赖 A 的路径保持 Y 为 1：

\begin{lstlisting}
Y_safe = (A AND B) OR ((NOT A) AND C) OR (B AND C)
\end{lstlisting}

这不是为了改变稳定真值表。因为当 \code{B AND C} 为 1 时，原表达式在 A=0 或 A=1 下本来也为 1；新增项只是在过渡期间提供连续路径。代价是多一个门和更多负载。是否值得加入，取决于这个输出是否会直接驱动危险控制线，或者是否会在毛刺窗口内被采样。

危险信号需要分级处理。并非所有组合输出都必须做成完全无毛刺：普通数据位通常可以容忍短暂过渡，只要后级在稳定的时钟边界采样；状态指示灯容忍度更高，短暂闪动人眼看不见，主要关心驱动能力；写使能和片选容忍度低，应避免异步毛刺，必要时寄存后再使用；复位、停机和电机使能容忍度极低，需要安全默认值、路径冗余和严格的时序检查。对这些信号，常见策略包括：把组合结果先送入触发器，在时钟边界后再使用；避免用毛刺风险高的组合式直接产生异步控制；必要时加入共识项或专用使能门；对外部执行机构再增加独立保护链路。

组合路径还需要时间预算。假设某个输出要在 20ns 内稳定，而路径包含输入缓冲、两级门、一级选择器和输出缓冲，就要把最坏延迟相加。下面的数字只是说明这种预算方法，不代表某个工艺或芯片规格。

\begin{longtable}{L{0.24\linewidth}L{0.26\linewidth}L{0.38\linewidth}}
\toprule
路径段 & 最坏延迟 & 说明 \\
\midrule
输入缓冲 & 2ns & 外部信号进入内部逻辑 \\
第一层门 & 3ns & 例如报警 OR 或安全条件 AND \\
第二层门 & 4ns & 例如 no\_fault 与请求合并 \\
选择器或缓冲 & 5ns & 输出选择或驱动增强 \\
布线和负载 & 3ns & 走线、电容、扇出影响 \\
合计 & 17ns & 若要求 20ns 内稳定，则余量 3ns \\
\bottomrule
\end{longtable}

负载也应纳入同一份时间预算。假设一个输出驱动 8 个 CMOS 输入，每个输入电容按 5pF 估算，走线和封装再贡献约 20pF，则总负载约为 60pF。若输出级等效电阻按 100$\Omega$ 粗估，时间常数约为 6ns；若后级要求在 10ns 内越过阈值，这条连接已经没有多少余量。若通过缓冲把 8 个后级分成两组，每个缓冲只承担较小负载，边沿和电平余量通常更容易满足。这个例子不替代仿真和数据手册，但它说明扇出不是“能接几个门”的静态数字，而是电流、电容、边沿和采样时刻共同形成的约束。

\begin{longtable}{L{0.2\linewidth}L{0.28\linewidth}L{0.42\linewidth}}
\toprule
估算项 & 数值 & 结论 \\
\midrule
输入电容 & $8 \times 5pF$ & 后级越多，总电容越大 \\
走线和封装 & 约 20pF & 板级和封装不是零负载 \\
总负载 & 约 60pF & 边沿速度需要重新估算 \\
等效 RC & $100\Omega \times 60pF \approx 6ns$ & 采样窗口紧时需要缓冲或分层驱动 \\
\bottomrule
\end{longtable}

\topic{扇出估算：每个后级输入都是一只小电容}
上表把负载当成一个给定的 60pF，这个数字从哪里来？答案是一条经验估算规则：每个 CMOS 输入贡献约 3–5pF（栅电容加引脚封装），每厘米走线对地平面再贡献约 1pF。一个输出驱动 10 个输入、走线 15cm，负载约为 $10\times4\mathrm{pF}+15\times1\mathrm{pF}\approx55$pF。驱动这样一个输出，本质上就是给这只总电容充电和放电，扇出问题的全部后果都从这一只电容展开。

边沿快慢可以用一阶 RC 粗估。从 74HC 数据手册的输出压降可以反推输出级等效电阻：额定 4mA 负载下 $V_\mathrm{OL}$ 约 0.2V（典型值，0.33V 为最坏值），对应约 50$\Omega$。于是时间常数 $\tau = R\cdot C = 50\Omega\times55\mathrm{pF}\approx2.8$ns，10\%–90\% 边沿约为 $2.2\tau\approx6$ns。若扇出加倍到 20 个输入、走线 20cm，负载约 100pF，边沿拖到约 11ns——已经与一级门自身的传播延迟相当，等于白白多串了一级门，而这级“门”不贡献任何逻辑功能。

\begin{longtable}{L{0.2\linewidth}L{0.36\linewidth}L{0.34\linewidth}}
\toprule
估算项 & 规则与算式 & 结果 \\
\midrule
后级输入电容 & 每输入 3–5pF，$10\times4$pF & 约 40pF \\
走线电容 & 约 1pF/cm，$15\mathrm{cm}\times1$pF & 约 15pF \\
合计负载 & $40+15$ & 约 55pF \\
边沿时间 & $2.2\times50\Omega\times55$pF & 约 6ns \\
扇出加倍后 & $2.2\times50\Omega\times100$pF & 约 11ns，堪比一级门延迟 \\
\bottomrule
\end{longtable}

高扇出还直接推高功耗。每次翻转搬运的电荷是 $Q=C\cdot V$，动态功耗为 $P=C\cdot V^{2}\cdot f$：55pF 在 5V、1MHz 下约 1.4mW，一块板上有几十根这样的重负载信号线，就是几十毫瓦白白发热。更隐蔽的是第二重代价：边沿变慢以后，后级输入在阈值附近停留的时间变长，接收门内部上下管同时微导通的窗口随之拉长，产生额外的短路功耗——慢边沿不仅自己慢，还让接收它的每一扇门更耗电。

缓冲器解决的正是这个物理问题。它不改变逻辑值，却把一个重负载拆成几个轻负载：先由普通输出驱动缓冲器，再由缓冲器分组驱动各片后级，每条分支电容小、边沿快，局部延迟反而更可控。代价是缓冲器自身的传播延迟要计入路径，所以缓冲层级不是越多越好——这又是一次功能、负载与延迟之间的取舍，与“逻辑上保持原值”并不矛盾。

若后续系统在 15ns 时就使用这个输出，稳定真值表仍然不够；输出可能还在传播或毛刺阶段。触发器和时钟边界的作用正是规定“什么时候保存”。组合逻辑在边界之间准备结果，时钟边界只应采样已经稳定、满足建立时间要求的信号。

时间预算要使用最坏情况，而不是典型情况。最小延迟说明某条路径最快可能多久变化，用来检查保持时间、窄脉冲和过早到达；典型延迟是常见条件下的大致延迟，便于估算，但不能作为设计边界；最大延迟是不利条件下最晚多久稳定，用来检查时钟周期和采样时刻是否足够；偏斜是不同路径或不同信号到达时间的差，用来解释毛刺、竞争和采样风险。若一条路径在实验室常温下 12ns 稳定，并不能证明它在低电压、高温、最大负载下也能在 15ns 内稳定。正式规格通常同时给出最小、典型、最大值；做时序分析时应以满足系统要求的最坏组合为准。

下面是几个基本时序词汇。建立时间指数据在采样时钟边沿到来前必须提前稳定的最短时间——结果不能刚到边沿才变化；保持时间指数据在采样时钟边沿之后仍必须继续稳定的最短时间——结果不能过早被下一条路径改写；时钟到输出指触发器被时钟采样后输出变化所需的时间，它是新数据进入下一段组合逻辑的起点；时钟偏斜指不同触发器看到同一时钟边沿的时间差，它会压缩某些路径的有效时序余量。

最基本的时序判断是：组合路径的最大延迟不能侵占建立时间，组合路径的最小延迟不能破坏保持时间，时钟偏斜不能被假定为零。换言之，组合逻辑不是“算完就行”，而是必须在规定的时钟边界之前算完，并且在边界附近保持稳定。

\topic{传播延迟沿路径累加：一条五级门的核算}
前面时间预算表里的数字是示意性的，现在用真实器件把同一条账重算一遍。回到安全启动控制板：从板级连接器到 \code{allow_start} 的最长路径，经过输入调理反相（一片 74HC04）和四级 NAND（两片 74HC00 依次产生 \code{n1}、\code{safe_chain}、\code{n3}、\code{allow_start}），共五级门。链式结构有一个简单却关键的性质：前一级的输出越过阈值后，后一级的延迟才开始计时，所以总延迟近似等于各级传播延迟之和。

查 74HC 系列数据手册（4.5V–5V 供电、50pF 负载电容、25°C 附近的标称条件），这一档门的单级传播延迟典型值约 9ns；手册同时给出保证最大值约 20ns，覆盖工艺离散、全温度范围和电压波动的最不利组合。不同厂商、不同温度档次的具体数字略有出入，正式设计必须以手头器件的手册为准，这里取代表值是为了说明算法。

\begin{longtable}{L{0.1\linewidth}L{0.44\linewidth}L{0.14\linewidth}L{0.14\linewidth}}
\toprule
路径级 & 器件与作用 & 典型 & 最大 \\
\midrule
第 1 级 & 74HC04，输入调理反相 & 9ns & 20ns \\
第 2 级 & 74HC00，产生 \code{n1} & 9ns & 20ns \\
第 3 级 & 74HC00，自反相得 \code{safe_chain} & 9ns & 20ns \\
第 4 级 & 74HC00，汇合得 \code{n3} & 9ns & 20ns \\
第 5 级 & 74HC00，自反相得 \code{allow_start} & 9ns & 20ns \\
合计 & 五级串联 & 45ns & 100ns \\
\bottomrule
\end{longtable}

两个合计数回答两个不同的问题。45ns 回答“台架上大概会看到什么”：实验台恰好常温、额定 5V、轻载，手头的芯片又只是少量样本，所以示波器上的实测值通常落在典型值附近——测量常见典型值，是因为实验室条件恰好全部偏有利。100ns 回答“设计必须容忍什么”：任何一片合格芯片都允许慢到最大值，换一批料、盛夏高温、电源跌到下限、负载加重，都可能把延迟推向 100ns。若采样边界按典型值设计，比如 60ns 就读取，台架上一切正常，量产后最不利的那批板子却会读进旧值。这就是为什么建立时间一类的检查必须用最大值核算，而典型值只能用来估量级。

还有两条细则。第一，数据手册把上升翻转（$t_\mathrm{PLH}$）和下降翻转（$t_\mathrm{PHL}$）分开给出，两者一般不相等，链式相加时每一级都应取两者中较大的那个。第二，50pF 是手册规定的测试条件；若实际负载更重，每级还要按扇出估算中的 RC 法则追加延迟——门的固有延迟和负载造成的延迟从来不是两个无关的量。

现在可以更系统地回答“现代芯片为什么延迟低”。第一层原因来自器件和工艺。更小的晶体管通常意味着更小的栅电容、更短的沟道、更短的局部连接和更低的单个门负载；同样的逻辑函数，在更小、更近的器件中传播，很多局部路径的绝对时间会下降。但这不是无条件成立。长连线、全芯片时钟、供电网络和散热会逐渐成为新的限制，所以现代芯片必须同时优化\hwkey{晶体管}和\hwtiming{连线}。

第二层原因来自片上集成。8086 时代，处理器与内存、外设和许多支撑逻辑之间存在大量片外通信；片外总线距离长、电容大、引脚数量有限，等待代价高。8086 已经用总线接口单元和执行单元的分工，以及指令预取队列，尝试把取指和执行重叠起来。到 80386，更多地址转换、保护和控制逻辑进入处理器内部，32-bit 数据通路和更丰富的片上控制减少了对外部协作的依赖。再往后，cache、浮点单元、内存控制器、图形或 I/O 控制器逐步进入同一颗芯片或同一封装，许多原本必须跨板级总线的路径被缩短为片上路径。

第三层原因来自存储层次。主存容量大，但距离核心远、访问路径长；寄存器和 L1 cache 容量小，却离执行单元近。现代芯片通过多级 cache 把最常用的数据和指令放在更近的位置。一次 L1 命中不是因为 DRAM 变得同样快，而是因为本次访问根本没有走到 DRAM。延迟降低来自“常见情况走短路径”，不是“所有路径都同样短”。

第四层原因来自流水线和并行。把一条很长的组合路径拆成多个较短阶段，每个阶段之间放入寄存器，就能提高时钟频率。这样做降低了每个时钟周期需要容纳的组合延迟，但并不等于每条指令的总周期数一定减少。现代处理器还会用分支预测、乱序执行、多个执行单元和内存预取，把将来可能需要的工作提前准备，把等待主存或分支结果的空洞尽量填满。它们降低的是常见执行路径上的可见等待。

因此，现代芯片的低延迟是一组工程选择的结果。器件缩小降低局部门延迟；门级库、缓冲和布局缩短关键路径；寄存器边界把长组合逻辑切成阶段；cache 把常用数据放近；预测和乱序把等待隐藏；片上集成减少跨芯片通信。若某次访问落在 \hwrisk{DRAM}、分支预测失败、cache 未命中或跨核心通信，延迟仍会显著上升。严肃的硬件分析必须同时说明\hwevidence{最快常见路径}和\hwrisk{最坏慢路径}。这些机制在本书后续章节（ISA、流水线、cache 与虚拟内存）和经典芯片与平台案例里会逐个展开成可检验的结构；数字硬件基础部分只需要先建立判断框架：低延迟来自让常见情况少走远路。

\begin{chipcase}
现代 CPU 的一个典型低延迟路径是：操作数已经在寄存器或 L1 cache 中，分支预测正确，所需执行单元空闲，结果能在相邻流水线阶段之间传递。此时，芯片用局部电路、短路径和寄存器边界把等待压到很低。相反，若数据不在 cache 中、分支预测错误或需要跨核心同步，同一条指令流会进入慢路径。现代芯片的工程目标不是让所有路径同样快，而是让高频路径短、让慢路径少发生，并在慢路径发生时尽量隐藏等待。
\end{chipcase}

还要区分\hwkey{延迟}和\hwtiming{吞吐率}。流水线可以让每个周期都有不同指令处在不同阶段，从而提高吞吐率；但一条具体指令从进入流水线到完成，仍然要经过多个阶段。多执行单元可以同时处理多条互不依赖的操作，但无法让一条有严格依赖链的操作跳过必要计算。cache 可以显著降低常见访问延迟，但不能让未命中的访问变成命中。理解现代芯片时，必须同时问两个问题：单次操作最快多久返回，连续大量操作每单位时间能完成多少。

还要区分反相器和缓冲器。反相器改变逻辑含义，缓冲器通常保持逻辑含义但增强驱动能力或隔离负载。很多时候，插入缓冲器不是为了“多一个门”，而是为了让一个信号能推得动更长的线或更多的后级。软件读者容易觉得保持值不变的缓冲器没有意义，但在硬件里，驱动能力本身就是功能的一部分。

最后还要检查扇入。扇出关心一个输出驱动多少后级，扇入关心一个门接收多少输入。四输入、八输入甚至更多输入的门在逻辑表达式里写起来很自然，但在硬件实现中可能被拆成多级结构；输入越多，门内部路径、电容和延迟也可能越大。把 \code{allow_start} 写成一个四输入 AND 并不意味着实际芯片里一定存在一个理想四输入 AND。实现可能把它拆成树形结构，也可能先合成若干中间信号。拆分方式会改变最长路径和毛刺形态。

\begin{lstlisting}
flat idea:
  allow_start = A AND B AND C AND D

tree implementation:
  left  = A AND B
  right = C AND D
  allow_start = left AND right
\end{lstlisting}

树形实现通常比线性串接更平衡，但会引入中间节点和不同路径长度。若某一路输入很晚才稳定，把它放在最后一级汇合可能有利于减少无效切换；若 A、B 来自同一连接器，先合成 \code{left} 可能便于局部布线。这里没有脱离功能的“最佳写法”，只有在功能、负载、延迟、布局和安全策略之间作出的实现选择。

对组合逻辑的最终检查，不应只停留在“真值表正确”。功能上，稳定输入下输出要符合真值表，否则是逻辑行为错误；默认状态上，悬空、断线、复位前的输入解释要有定义，否则异常时会误启动或误使能；驱动能力上，每个输出要能带动实际负载，否则电平被拖低、边沿过慢；路径延迟上，最长组合路径要能在后级使用前稳定，否则后级采到旧值或过渡值；毛刺风险上，多路径汇合不能产生危险窄脉冲，否则写使能、复位、片选会被误触发；低有效信号上，信号命名和真值表要一致，否则使能和禁止方向会接反；物理边界上，长线、外部输入和按钮抖动都要处理，否则干扰和抖动会进入逻辑核心。

对安全启动控制板，可以把这些检查整理成矩阵。矩阵的价值在于把“正常输入”“边界输入”“异常输入”和“时间条件”放在一起，防止只验证最容易通过的几行真值表。

\begin{longtable}{L{0.18\linewidth}L{0.34\linewidth}L{0.38\linewidth}}
\toprule
检查类别 & 要施加的条件 & 期望结果 \\
\midrule
正常允许 & 请求有效、盖闭合、急停正常、无故障 & \code{allow_start}=1，执行命令等待时序边界 \\
正常禁止 & 任一安全条件为 0 & \code{allow_start}=0 \\
故障禁止 & 温度或电流报警为 1 & \code{fault_lamp}=1，\code{allow_start}=0 \\
断线默认 & 保护盖、急停或报警链路断开 & 进入禁止或报警一侧 \\
上电过程 & 电源未稳或复位保持 & \code{motor_enable_cmd}=0 \\
按钮抖动 & 原始按钮快速断续 & 内部请求不产生多次启动事件 \\
最坏负载 & 故障灯、后级逻辑和连接器同时连接 & 电平仍满足后级阈值 \\
最坏时序 & 输入在不利延迟组合下变化 & 后级只在稳定后采样或使用 \\
\bottomrule
\end{longtable}

常见失败也应明确记录。许多硬件问题并不是因为设计者不知道 AND、OR、NOT，而是因为把抽象层次接错：悬空输入让系统偶尔误启动或误报警，根源是没有默认上拉/下拉或输入调理；低有效误读让复位、片选或急停方向相反，根源是信号约定未先于公式建立；电平不兼容让单独测试正常的电路接入后失效，根源是前后级阈值和驱动能力不匹配；毛刺误触发让最终值正确的电路出现短促动作，根源是危险输出直接使用了组合信号；扇出过大让边沿变慢、高电平不够高，根源是一个输出承担了过多负载；时序乐观让常温可用的电路在高温或低压下失效，根源是用典型值代替最坏值做核算；无关项误用让异常状态进入允许输出，根源是把可能发生的故障组合当作无关项。

组合电路需要两种同时成立的视角。第一，抽象视角：逻辑门实现布尔函数，组合网络实现加法、选择、译码和比较。第二，物理实现视角：每个输出都来自具体电流路径，每条路径都有延迟、负载和阈值，每个危险输出都需要明确的使用时刻。缺少前者，无法构造大系统；缺少后者，系统只是在纸上正确。

本章结束时，读者应能整理出一份小型硬件设计包。它不要求购买昂贵仪器，也不要求完成工业级认证；它要求把抽象、实现和验证记录放在同一份文档里：一张信号表，列出原始输入、内部语义信号、组合输出和危险输出及其有效电平；一张电平表，列出前级 $V_\mathrm{OH}/V_\mathrm{OL}$、后级 $V_\mathrm{IH}/V_\mathrm{IL}$、噪声余量和负载条件，全部用最坏值；一张路径表，写出关键门在每行输入下的上拉、下拉、悬空和冲突状态；一张实现表，写明使用的门芯片、缓冲器、上拉/下拉、保护和输出驱动；一张时序表，写明最长路径、最小路径、毛刺风险、后级使用时刻和危险输出策略；一张测试表，写明施加条件、期望结果、测量工具、测点和波形记录；外加几张案例卡，把 7400/4000/8086/80386/现代 CPU 与本章概念真正关联起来。若后续章节要把这个组合逻辑接入寄存器、CPU 或总线，这份设计包就是接口约定。

对安全启动控制板的检查也应按执行顺序展开，而不是只看最终结果。第一步在断电和上电爬升阶段确认 \code{motor_enable_cmd} 被保持为 0；第二步测量每个原始输入在松开、按下、断线和报警状态下的默认电平；第三步用示波器确认按钮抖动、传感器边沿和故障输入不会在阈值附近长期停留；第四步用逻辑分析仪或等效测试记录内部语义信号是否只产生规定事件；第五步施加正常允许、单故障、双故障和断线条件，确认组合输出满足真值表；第六步在最坏负载和最坏时序假设下确认后级只在稳定后使用结果。

实验记录应当能让第三方复查。只写“测试通过”不够；至少要记录条件、预期结果、测量数据和结论。下面的模板适合本章的小组合模块，也适合后续寄存器和总线实验继续扩展。

\begin{longtable}{L{0.13\linewidth}L{0.17\linewidth}L{0.2\linewidth}L{0.2\linewidth}L{0.1\linewidth}}
\toprule
测试项 & 施加条件 & 期望结果 & 实测电压或波形 & 结论 \\
\midrule
上电默认 & 电源爬升、复位保持 & \code{motor_enable_cmd}=0 & 记录关键节点电压与复位释放时间 & 合格或失败 \\
按钮输入 & 松开、按下、抖动、断线 & 内部请求只产生规定事件 & 记录阈值穿越和消抖后信号 & 合格或失败 \\
故障禁止 & 温度或电流报警有效 & \code{fault_lamp}=1，\code{allow_start}=0 & 记录报警输入与组合输出延迟 & 合格或失败 \\
最坏负载 & 连接全部后级和指示灯驱动 & 电平余量和边沿仍满足规格 & 记录高低电平、上升/下降时间 & 合格或失败 \\
毛刺检查 & 输入按不利时序翻转 & 危险输出不响应窄脉冲 & 记录示波器或逻辑分析仪波形 & 合格或失败 \\
\bottomrule
\end{longtable}

若使用逻辑分析仪，还应记录采样率、触发条件和被测信号名称；若使用示波器，还应记录探头倍率、带宽限制、地线接法和触发边沿。测量工具本身也会影响电路：长地线可能引入额外振铃，高电容探头可能拖慢边沿。因此，实验记录不仅证明电路行为，也证明测量方法没有把问题掩盖或放大。

实际测量时，探头应尽量接在接收端附近，而不只接在驱动端。驱动端波形漂亮，不代表穿过连接器、长线和负载后的接收端仍有足够余量。示波器观察快速边沿时应优先使用短地弹簧或低感抗接地方式，避免探头地线把额外振铃带入测量结果；逻辑分析仪只给出阈值后的数字结果，不能证明边沿是否干净；万用表只能证明低速或静态电压，不能证明毛刺和建立时间。正式记录应注明测点、探头、阈值、触发条件和环境条件——只证明常温空载，不等于证明最坏条件。

异步输入进入时钟边界时可能出现亚稳态。亚稳态指触发器在采样边沿附近遇到变化中的输入，内部节点可能暂时停在不能立刻判定为 0 或 1 的状态。边界原则是：外部按钮、传感器、中断线和跨时钟域信号，不应未经同步就直接控制内部状态或危险输出。组合逻辑可以整理当前事实，但“何时保存”必须由时序电路给出约定。

\topic{把电平特性从“能亮”推进到“有裕量”}
真正做硬件时，还必须把“能看到 0/1”推进到“在误差、噪声、负载和温度变化下仍然可靠”。一个输入脚不只是在某个理想电压上突然改变含义，而是由输入低阈值、输入高阈值、输出低电平、输出高电平共同定义可用范围。输出端能给出的高电平必须高于下一级认可的高阈值，输出端能拉到的低电平必须低于下一级认可的低阈值，中间差额就是噪声余量。

\begin{longtable}{L{0.18\linewidth}L{0.38\linewidth}L{0.34\linewidth}}
\toprule
参数 & 含义 & 设计时要回答的问题 \\
\midrule
\code{VOH} & 输出为 1 时保证达到的最低高电平 & 在最大负载下是否仍高于下一级 \code{VIH} \\
\code{VOL} & 输出为 0 时保证不超过的最高低电平 & 在灌电流时是否仍低于下一级 \code{VIL} \\
\code{VIH} & 输入被识别为 1 的最低电压 & 外部信号、上拉和噪声是否能跨过阈值 \\
\code{VIL} & 输入被识别为 0 的最高电压 & 慢边沿和串扰是否会停在不确定区 \\
噪声余量 & 输出保证值到输入阈值之间的余量 & 板级噪声、地弹和电源纹波是否小于余量 \\
\bottomrule
\end{longtable}

这张表比“LED 亮不亮”更接近工程真实。LED 亮只能说明某个输出能驱动肉眼可见负载，不能证明它满足数字输入阈值；示波器上看到高电平接近电源，也不等于在最大扇出、最坏温度和最长连线下仍可靠。每个数字信号都要同时看电平范围、驱动能力、边沿速度和负载。

\begin{lstlisting}
noise margin check:
  high_margin = VOH_min - VIH_min
  low_margin  = VIL_max - VOL_max
  pass only if expected_noise < both margins
\end{lstlisting}

\topic{CMOS 反相器的关键不是“一个管上、一个管下”}
CMOS 反相器常被画成上拉 PMOS 和下拉 NMOS。这个图很有用，但只说“输入 0 时上管开、下管关”还不够。实际电路要处理三个非理想问题：输出节点有电容——栅电容、扩散电容、连线电容和下级输入电容累加，翻转需要充放电时间，这就是传播延迟的来源；输入在阈值附近缓慢变化时，上下管可能同时部分导通形成短路电流，慢边沿因此既增加功耗也增加噪声；驱动多个输入或长线时等效电容更大，边沿更慢，延迟和裕量都被压缩。同一张清单上还有两条板级事实：多个输出同时翻转产生的瞬态电流会让地弹或电源跌落，输入阈值附近可能被误触发；高阻输入容易被漏电和耦合影响，必须用上拉、下拉或明确驱动源给出归宿。

因此，数字硬件基础部分的“开关”要升级为“受控电流路径”。当输出从 0 变 1 时，上拉网络不是瞬间把节点变成电源，而是向负载电容充电；当输出从 1 变 0 时，下拉网络把负载电容放电。电容越大、电阻越大，边沿越慢。

\topic{扇出、上拉和开漏输出要按电流读}
许多初学错误来自只看电压不看电流。一个输出脚能否驱动多个输入，要看它在保持合法 \code{VOH/VOL} 时能提供或吸收多少电流。TTL、CMOS、开漏、三态输出的驱动方式不同，不能只按“都是 5V 逻辑”混接。

\begin{longtable}{L{0.18\linewidth}L{0.38\linewidth}L{0.34\linewidth}}
\toprule
输出形式 & 电流路径 & 常见用途和风险 \\
\midrule
推挽输出 & 上拉和下拉都由有源器件驱动 & 边沿快，但两个推挽输出不能直接并联 \\
开漏/开集 & 只能主动拉低，高电平靠上拉电阻 & 适合线与、总线仲裁，但上升沿较慢 \\
三态输出 & 可驱动高、低或断开 & 总线共享，但必须保证同一时刻唯一驱动者 \\
弱上拉/下拉 & 用较大电阻给默认电平 & 适合防悬空，不适合驱动大负载 \\
\bottomrule
\end{longtable}

开漏输出尤其适合解释“电压是结果，电流路径才是原因”。当任意设备拉低总线时，电流从电源经上拉电阻流入该设备到地，总线电压变为低；当所有设备都断开时，上拉电阻慢慢给总线电容充电，电压回到高。上拉电阻太大，上升沿慢；太小，拉低时电流大。I2C、按键输入和许多中断线都能用这个模型理解。

\begin{lstlisting}
open-drain line:
  any driver pulls low  -> line = 0
  all drivers released  -> pull-up charges line to 1
  rise time depends on pull-up R and bus capacitance C
\end{lstlisting}

\topic{去耦电容是本地瞬态电源}
门翻转的电流从哪里来？不是从稳压器“立刻”流过来——稳压器的响应以微秒计，而门的边沿以纳秒计。输出翻转时对负载电容充放电的电荷，加上内部上下管瞬间同时微导通的贯穿电流，都只能先从电源脚附近的储能元件取得。如果附近没有储能，电荷只能沿电源走线从远处赶来，而走线不是理想导线：每厘米约贡献 5–10nH 电感（含回流路径），电感上的电压满足 $v=L\cdot di/dt$，电流变化越快，走线上感应出的压降越大。

代入一组典型数字。8 路输出同时翻转，负载各 50pF，搬运的总电荷为 $Q=8\times50\mathrm{pF}\times5\mathrm{V}=2$nC；在 5ns 内完成，意味着瞬态电流约 $2\mathrm{nC}/5\mathrm{ns}=0.4$A。若这 0.4A 要穿过约 20nH 的走线电感，感应电压为 $20\mathrm{nH}\times0.4\mathrm{A}/5\mathrm{ns}=1.6$V——电源脚瞬时跌落、地脚瞬时抬高（后者俗称地弹），整片芯片看到的“5V”在剧烈抖动，输入阈值跟着晃，邻近信号可能被误判。地弹不是器件坏了，而是瞬态电流与走线电感的必然乘积。

每颗芯片就近放一只 0.1$\mu$F 瓷片电容，就是给瞬态电流修一座本地水库：尖峰电流先由它供给，走线电感只需在事后把电荷慢慢补回。同样 2nC 改由 0.1$\mu$F 供给，电压跌落只有 $\Delta V=Q/C=2\mathrm{nC}/0.1\mu\mathrm{F}=20$mV，比 1.6V 低约两个数量级。“贴近电源脚”是电气要求而非工艺美观：电容与芯片之间每多一毫米走线，水库与芯片之间就多串一段电感，本地储能的效果就打一分折扣。所以布局规则是紧贴电源脚、短焊盘、短回流路径。

0.1$\mu$F 只管得住快而小的口袋，整板还需要分层：电源入口的 10–100$\mu$F 大电容应付更大、更慢的电流波动（继电器吸合、数码管扫描、背光开启），稳压器负责直流供给。按时间尺度分层供电：快的层级小而近，慢的层级大而远，各自接住上一层来不及响应的那一段。

\topic{走线多短才算“短”：信号完整性初步}
信号在走线上的传播速度是有限的：FR4 板材上约为 15cm/ns，大约是光速的一半。只要走线延迟远小于信号边沿时间，线上各点就来不及显现分布特性，可以把整条走线当作一只集中电容——这就是“短”线。常用的经验法则是：走线延迟小于边沿时间的六分之一（有的教材取四分之一），即可按集中参数处理。前面的扇出估算把走线只算作每厘米 1pF，依据正是这个近似。

74HC 的边沿约 5–8ns，代入法则：临界长度 $\approx5\mathrm{ns}\times15\mathrm{cm/ns}\div6\approx12$cm。本章的面包板和实验连线都在这个量级之下，所以把走线当集中电容是站得住的——本书低速实验大多安全地待在“短”线区。反过来，一旦边沿变快（现代器件可低于 1ns）或距离变长（跨板、背板、电缆），短线近似就会失效，同一根线必须改按传输线看待。

长线上的核心现象是反射。阻抗不连续的地方，行波会部分折返：无端接的长线末端近似开路，入射波几乎全反射回来，与入射波叠加后，接收端出现过冲、振铃甚至回沟，一个边沿可能被看成多次穿越阈值——接收端误以为信号翻转了好几回。端接的思路是为行波安排能量归宿：末端并联电阻到地或电源，或源端串联电阻，使反射系数接近零，波形一次到位。反射不是噪声玄学，它的幅度和极性都可以从阻抗差算出。

串扰则来自相邻走线之间的互容与互感。两条长平行线中，一条快速翻转会在另一条上感应出窄毛刺；平行段越长、间距越近、边沿越快，耦合越强。处理办法是拉开间距、缩短平行段、在敏感信号之间夹一根地线。这些现象在 HC 实验里只会以“边沿有点圆、偶尔带振铃”的形式露面，但同一条物理规律决定了高速主板为什么讲究受控阻抗与端接电阻。

\topic{毛刺不是玄学，而是路径延迟不一致}
组合逻辑里不同输入路径有不同延迟。当多个输入几乎同时变化时，中间节点可能短暂经过错误组合，输出形成窄脉冲，这就是毛刺。若毛刺只被后级组合逻辑看到，可能不影响最终稳定值；若毛刺被锁存器、计数器、异步复位或外部设备采样，就会变成真实错误。

毛刺最常见的来源有四种。重汇聚路径：同一输入经不同延迟路径再汇合，处理办法是改写逻辑、增加共识项，或让后级只在时钟边界采样。译码器切换：地址多位同时变化，片选可能短暂重叠或空洞，处理办法是地址先锁存、片选加使能窗口。机械按键：触点弹跳导致多次高低变化，处理办法是 RC、施密特、同步消抖或状态机消抖。异步输入：外部事件不按本机时钟变化，处理办法是同步器和边沿检测分开设计。

组合逻辑允许在稳定前经历暂态，时序元件只应在规定窗口采样稳定结果。若把组合毛刺直接接到计数器时钟、异步清零或设备写使能，系统会表现得像随机错误。

\topic{数字硬件基础部分实验应记录波形而不只记录现象}
一块面包板、一个反相器、一只按键和一台示波器，就能做出高质量的数字硬件基础部分实验。目标不是炫耀器件数量，而是把电平、阈值、边沿、负载和毛刺变成可以直接观察的量。

\begin{longtable}{L{0.18\linewidth}L{0.42\linewidth}L{0.3\linewidth}}
\toprule
实验 & 操作 & 必须记录 \\
\midrule
输入阈值扫描 & 用电位器缓慢改变输入电压 & 输出翻转点、滞回是否存在 \\
负载电容实验 & 在输出加不同电容负载 & 上升/下降时间和传播延迟变化 \\
上拉电阻实验 & 改变开漏线的上拉电阻 & 上升沿、低电平电流和噪声余量 \\
按键抖动观察 & 直接观察按键输入波形 & 抖动持续时间和次数 \\
译码毛刺观察 & 让多位地址同时变化 & 片选是否短暂重叠或误触发 \\
\bottomrule
\end{longtable}

这些实验报告应写出测量条件：电源电压、器件型号、负载、电阻电容值、示波器探头倍率、时基和触发方式。没有测量条件的波形很难复现，也很难判断结论是否能迁移到另一块板。

\section*{章末练习}

\begin{longtable}{L{0.18\linewidth}L{0.39\linewidth}L{0.33\linewidth}}
\toprule
任务 & 要完成的产物 & 检查要点 \\
\midrule
按钮输入 & 为 \code{start_button_raw} 写出上拉或下拉方案、有效电平、断线默认值和消抖策略 & 松开、按下、抖动、断线四种状态都有明确内部语义 \\
电平约定 & 给出一组前级 $V_\mathrm{OH}/V_\mathrm{OL}$ 与后级 $V_\mathrm{IH}/V_\mathrm{IL}$，计算高低噪声余量 & 说明余量为正、为零、为负时分别意味着什么 \\
路径表 & 任选 NAND 或 NOR，列出四行真值表，并为每一行写出上拉与下拉路径状态 & 不出现稳定输入下的悬空或强冲突 \\
低有效信号 & 将 \code{reset_n}、\code{chip_select_n} 或急停链路转换为正逻辑语义信号 & 说明原始电平、内部语义和默认安全值之间的关系 \\
故障汇总 & 为 \code{temp_alarm} 与 \code{current_alarm} 设计故障灯和故障码逻辑 & 单故障、双故障、无故障和非法输入都有定义 \\
1-bit ALU & 写出 AND、OR、XOR、ADD、SUB 的候选结果、\code{op}、\code{B_invert} 和最低位 \code{Cin} 规则 & 区分逐位逻辑路径、加减法进位路径和标志解释 \\
NAND-only 实现 & 把 \code{allow_start} 的安全表达式改写成二输入 NAND 网络，并命名所有反相中间信号 & 稳定真值表等价，同时说明门级层数、负载和毛刺风险 \\
芯片引脚映射 & 任选 74HC00 或 7400 类芯片，把一个逻辑式映射到电源脚、输入脚、输出脚和未用门处理 & 说明去耦、未用输入默认值、扇出和封装方向 \\
输入保护 & 为一个板外按钮或传感器写出限流、滤波、ESD/TVS、施密特输入或缓冲方案 & 说明保护不改变语义，但限制异常电压、电流和噪声 \\
五层边界图 & 任选 \code{start_button}、\code{cover_closed} 或 \code{temp_alarm}，画出外部物理量、保护边界、阈值整形、语义归一和组合逻辑五层 & 每层都要写出输入、输出、风险和验证方法，不能把原始引脚直接当作内部 bit \\
引脚等效模型 & 为一个输入和一个输出分别列出漏电、电容、阈值、驱动电流、传播延迟或绝对最大额定值 & 能解释为什么逻辑变量模型不足以判断扇出、长线、LED 驱动和跨电源域连接 \\
毛刺分析 & 分析 \code{Y=(A AND B) OR ((NOT A) AND C)} 在 B=1、C=1、A 翻转时的风险 & 能说明共识项为什么不改变稳定真值表，却改变过渡行为 \\
安全闭环 & 按“原始输入、输入调理、语义信号、组合判断、时序边界、输出驱动、验证记录”重写安全启动控制板规格 & 每一层都能追溯到上一层，危险输出不能直接依赖未经检查的组合结果 \\
芯片案例 & 分别用 7400、4000 CMOS、8086、80386 和现代 CPU 解释本章一个概念；6502、Z80、8051、80486、Pentium 作为选做，等学完经典芯片与平台案例后回看 & 每个案例都要指出对应的电平、路径、组合逻辑、片上集成或低延迟机制 \\
测量计划 & 为按钮输入或故障输出写出万用表、示波器、逻辑分析仪各自要证明的内容 & 不把静态电压测量误认为完整的时序和毛刺检查 \\
时序词汇 & 用一条组合路径解释建立时间、保持时间、时钟到输出和时钟偏斜 & 能说明最大延迟、最小延迟和采样边界的关系 \\
实验记录 & 按测试项、条件、期望结果、实测波形和结论整理一次安全启动控制板实验 & 记录足以让第三方复查，不只写“通过” \\
板级边界 & 说明参考地、回流路径、去耦、串扰或长线振铃如何影响某个输入 & 能把布尔错误和物理边界错误区分开 \\
快慢路径 & 为现代 CPU 写出一个快路径和一个慢路径案例 & 能说明低延迟来自常见路径短，而不是所有路径同样短 \\
上拉阻值 & 为一条 5V 开漏总线（总线电容 100pF，驱动器灌电流上限 4mA）选择上拉电阻，要求从低电平上升到 $V_\mathrm{IH}=3.5$V 的时间不超过 1$\mu$s & 同时算出阻值上界和下界，并验证所选标称值两头都满足 \\
三态总线互斥 & 为共享一条数据总线的两个三态驱动器写出使能逻辑，并分析选择信号切换瞬间的风险 & 任何时刻至多一个驱动器使能，切换要先断后通 \\
选择器实现函数 & 用 8 选 1 选择器实现 \code{F(A,B,C)}$=\Sigma m(1,3,5,6)$，再改用 4 选 1 选择器、以 A 作数据输入实现同一函数 & 每个数据输入的常量或函数都要与真值表逐行一致 \\
延迟链核算 & 一条 4 级 74HC 路径，每级典型 9ns、最大 22ns，要求信号在采样前 90ns 稳定 & 分别用典型值和最大值核算，说明哪一个才能作为设计依据 \\
去耦电容选择 & 为一块 8 路输出可能同时翻转的 74HC 小板选择去耦方案，估算本地电容上的电压跌落 & 说明就近原则与板级大电容的分工，跌落估算要有数值 \\
电平转换 & 3.3V CMOS 输出（最差 $V_\mathrm{OH}=3.0$V）驱动 4.5V 供电的 74HC 输入（$V_\mathrm{IH}=3.15$V），给出可行方案 & 先核算直接连接的余量，再给出器件级方案并重算余量 \\
\bottomrule
\end{longtable}

\section*{参考解答与检查要点}

本节给出参考答案。实际工程方案允许存在不同器件和参数选择，但结论必须有电平、路径、时序和测量结果支撑。若答案只给出布尔式而没有默认值、负载、毛刺或测量边界，不能视为完整解答。

\begin{longtable}{L{0.18\linewidth}L{0.44\linewidth}L{0.28\linewidth}}
\toprule
任务 & 参考解答 & 必须保留的检查点 \\
\midrule
按钮输入 & 可采用上拉输入：松开时 \code{start_button_raw}=1，按下接地为 0，再经反相、消抖和同步得到 \code{start_request}=1；也可采用下拉输入，关键是松开、按下、断线都要有定义。 & 不允许松开或断线时悬空；按钮抖动不能形成多次启动事件。 \\
电平约定 & 例如 $V_\mathrm{OH}=3.0V$、$V_\mathrm{IH}=2.1V$，高余量 0.9V；$V_\mathrm{OL}=0.25V$、$V_\mathrm{IL}=0.7V$，低余量 0.45V。余量为正表示最坏条件仍可识别，余量为零表示没有抗扰空间，余量为负表示电平约定失败。 & 使用最坏值，不能用典型值替代核算。 \\
路径表 & 以 NAND 为例，A=B=1 时下拉串联路径完整、输出为 0；其余三行至少一个下拉开关断开，上拉网络提供高电平。NOR 则对应并联下拉和串联上拉。 & 每一行都要说明上拉、下拉、悬空和冲突。 \\
低有效信号 & \code{reset_n}=0 表示复位有效，可转换为 \code{reset_active = NOT reset_n}；\code{chip_select_n}=0 表示选中，可转换为 \code{chip_selected = NOT chip_select_n}。 & 名称必须反映极性，默认安全状态要明确。 \\
故障汇总 & \code{fault_lamp = temp_alarm OR current_alarm}；\code{no_fault = NOT fault_lamp}；双故障可输出专用多故障码 \code{11}。 & 任一故障都禁止启动，非法或未知组合不应进入允许状态。 \\
1-bit ALU & \code{AND=A AND B}，\code{OR=A OR B}，\code{XOR=A XOR B}；ADD 使用 \code{B_invert=0,Cin=0}；SUB 使用 \code{B_invert=1,Cin=1}，即 \code{A+(NOT B)+1}。 & \code{op}、\code{B_invert}、最低位 \code{Cin} 必须一致，标志位要说明无符号或补码解释。 \\
NAND-only 实现 & 可写为 \code{n1=cover_closed NAND emergency_ok}、\code{safe_chain=n1 NAND n1}、\code{n2=start_request NAND no_fault}、\code{request_ok=n2 NAND n2}、\code{n3=safe_chain NAND request_ok}、\code{allow_start=n3 NAND n3}。 & 反相中间信号必须命名清楚，NAND-only 等价不代表延迟和毛刺也相同。 \\
芯片引脚映射 & 以 74HC00 为例，一路 NAND 接 \code{cover_closed/emergency_ok} 得到 \code{safe_chain_n}，第二路 NAND 两输入并接 \code{safe_chain_n} 得到 \code{safe_chain}；\code{VCC/GND} 接正确电源并就近去耦，未用 CMOS 输入接确定电平。 & 检查封装方向、引脚编号、未用门、扇出、输出驱动和电源去耦。 \\
输入保护 & 板外按钮可采用串联电阻、上拉/下拉、RC 滤波、TVS/ESD 保护和施密特输入；高速输入不能使用过大 RC。 & 保护电路限制异常电压和电流，但不能破坏有效边沿和响应时间。 \\
五层边界图 & 以 \code{start_button} 为例，\code{raw} 表示触点和线缆；保护层包含限流、ESD 和默认上拉/下拉；整形层包含 RC 或施密特输入；语义层把低有效、消抖和反相整理成 \code{start_request}；组合层只使用 \code{start_request}。 & 每层输出都要比上一层更接近可靠 bit，不能跳过保护和阈值整形环节。 \\
引脚等效模型 & 输入端可写成阈值比较器加输入电容、漏电和保护结构；输出端可写成受控上拉/下拉开关加等效电阻、电流限制和延迟。这个模型能解释上拉太弱、负载过重、长线振铃、输出并接和跨电源域损坏。 & 数据手册最坏值优先；不要用“仿真里是 1”替代引脚电气检查。 \\
毛刺分析 & 对 \code{Y=(A AND B) OR ((NOT A) AND C)}，当 B=C=1 且 A 翻转时，两条路径延迟不同，Y 可能短暂掉到 0；加入共识项 \code{B AND C} 可保持静态 1。 & 稳定真值表不变，但过渡行为改变；危险输出要寄存或专门处理。 \\
安全闭环 & 规格应写成：原始输入先经保护、默认值和消抖，得到语义信号；组合逻辑形成 \code{fault_lamp/no_fault/allow_start}；危险输出经时序边界和驱动级形成 \code{motor_enable_cmd}。 & 原始触点不能直接驱动执行机构；复位和断线默认必须安全。 \\
芯片案例 & 7400 展示门芯片电气约定；4000 CMOS 展示高输入阻抗和低静态功耗；6502/Z80 展示外部总线微处理器；8051 展示片上 I/O；8086 展示预取和复用总线；80386 展示 32-bit 和保护；80486/Pentium 展示 cache、流水线和预测。 & 每个案例都要绑定到本章概念，不能只列名称。 \\
测量计划 & 万用表确认静态高低电平；示波器观察抖动、边沿、振铃和毛刺；逻辑分析仪记录内部语义事件和时序关系。 & 三类工具不能互相替代，测点应靠近接收端。 \\
时序词汇 & 建立时间是采样沿前必须稳定的时间；保持时间是采样沿后必须继续稳定的时间；时钟到输出是触发器输出延迟；时钟偏斜是不同触发器看到时钟沿的差异。 & 最大延迟影响建立，最小延迟影响保持。 \\
实验记录 & 记录测试项、条件、期望结果、实测电压或波形、逻辑分析仪结果和结论；注明探头倍率、触发条件、采样率、电源电压和负载。 & 只写“测试通过”不算结论，记录必须能让第三方复查。 \\
板级边界 & 地弹会压缩噪声余量；回流路径不连续会增加环路面积；去耦不足会造成供电跌落；串扰和振铃会让接收端误判。 & 区分布尔函数错误和物理边界错误。 \\
快慢路径 & 快路径：操作数在寄存器或 L1，分支预测正确，旁路可用；慢路径：cache miss、TLB miss、分支错误、跨核心同步或 DRAM 访问。 & 低延迟来自常见路径短，不代表所有路径都短。 \\
上拉阻值 & 上界由上升时间决定：$t=R\cdot C\cdot\ln\frac{V_\mathrm{CC}}{V_\mathrm{CC}-V_\mathrm{IH}}$，代入得 $t\approx R\times100\mathrm{pF}\times1.20$，由 $t\le1\mu$s 得 $R\le8.3$k$\Omega$；下界由低电平电流决定：拉低时电阻电流 $(5-0.4)\mathrm{V}/R\le4$mA，得 $R\ge1.15$k$\Omega$。取标称 4.7k$\Omega$：上升时间约 $4.7\mathrm{k}\Omega\times100\mathrm{pF}\times1.20\approx0.57\mu$s，低电平电流约 $0.98$mA，两关都过。 & 只算上升时间会漏掉拉低电流这一关；$\ln$ 因子来自电容充电曲线，不是凑出来的系数。 \\
三态总线互斥 & 使能由同一选择信号经 one-hot（独热）译码产生，并在时序边界上先全部断开、下一拍再接通新驱动者（先断后通）。若用 \code{EN_B = NOT EN_A} 组合派生，反相器延迟窗口内会出现两者同通或同断；同通时两个推挽输出并联，可能一个拉高一个拉低，形成冲突电流。 & 使能信号属于危险输出，要按毛刺与时序边界管理；分析切换瞬间必须写出电流路径。 \\
选择器实现函数 & 8 选 1：选择端接 A、B、C，数据端 \code{D1=D3=D5=D6=1}，其余接 0。4 选 1：选择端接 B、C，逐行对照得 \code{D0=0}（m0、m4 均为 0）、\code{D1=1}（m1、m5 均为 1）、\code{D2=A}（m2=0、m6=1）、\code{D3=NOT A}（m3=1、m7=0）。 & 每个数据输入都要能追溯到真值表的两行；选择端变量顺序一旦接反，整张映射全错。 \\
延迟链核算 & 典型合计 $4\times9\mathrm{ns}=36$ns；最大合计 $4\times22\mathrm{ns}=88$ns，仍满足 90ns 的要求，但余量只有 2ns。若误用典型值，会误以为余量 54ns。设计依据只能是最大值：$88\mathrm{ns}\le90\mathrm{ns}$。 & 建立类检查用最大值；典型值只用于估量级，不能写进设计边界。 \\
去耦电容选择 & 8 路同时翻转搬运电荷 $Q=8\times50\mathrm{pF}\times5\mathrm{V}=2$nC；由就近的 0.1$\mu$F 供给，跌落 $\Delta V=Q/C=20$mV。若无本地电容，约 0.4A（2nC/5ns）的瞬态电流穿过约 20nH 走线电感，将引起约 $L\cdot di/dt=1.6$V 的扰动。方案：每片 0.1$\mu$F 瓷片紧贴电源脚，电源入口再放 10–47$\mu$F。 & 就近是电气要求而非工艺美观；只写“加 0.1$\mu$F”而没有位置与数值估算不算完整。 \\
电平转换 & 直接连接的高电平余量为 $3.0-3.15=-0.15$V，必然失败。方案之一：用一片 74HCT（4.5V 供电、TTL 输入阈值 $V_\mathrm{IH}=2.0$V）作第一级整形，新余量 $3.0-2.0=1.0$V，输出已是标准 CMOS 电平；也可用开漏输出加上拉到 4.5V，或专用电平转换芯片。反向（4.5V 驱动 3.3V）可用串联电阻加钳位二极管，或选用耐 5V 输入的器件。 & 先证明原方案余量为负，再谈转换；不能只写“中间加一级”而不重算新余量。 \\
\bottomrule
\end{longtable}

\begin{classicnote}
本章的掌握标准不是“会背 AND、OR、NOT、XOR 的真值表”，而是能把每个抽象还原为具体的硬件行为。

\begin{itemize}
\item 电平：前一级的最差高/低输出，必须落在后一级输入阈值允许的范围内，并留下噪声余量。
\item 路径：任何输出节点都要能写成“高电平路径是否导通、低电平路径是否导通、是否存在冲突电流、是否存在悬空”四项。
\item 门级：任选一个二输入门，列出 4 行真值表，再为每一行写出上拉网络和下拉网络状态。
\item 组合：半加器的 Sum 是 XOR，Carry 是 AND；全加器是在同一张真值表中纳入进位输入后的组合电路。
\item 延迟：同一个逻辑函数可以有不同门级网络，最长路径不同，稳定时间也不同。
\item 负载：一个门的输出不是无限共享变量，后级数量、线长和缓冲策略会改变电气行为。
\item 边界：外部输入、上电复位、开漏共享线、三态总线和电源域转换都必须先满足边界约定，才能进入普通组合逻辑。
\item 安全：内部允许判断不等同于执行机构命令，危险输出必须说明默认状态、使用时刻和毛刺处理策略。
\item 芯片案例：7400/4000 系列展示标准门芯片的电气约定；8086/80386 展示组合控制和数据通路如何扩展成处理器；现代 CPU 展示 cache、流水线、预测和片上集成如何降低常见路径的可见等待。
\item 检查题：若组合网络输出短暂从 0 跳到 1 又回到 0，而最终值正确，它是否一定构成设计错误？答案取决于后级是否在毛刺期间采样。
\end{itemize}

经典系统教材常把抽象的位函数和它的物理实现放在同一页上看。只看位函数会漏掉电平、驱动和时序；只看器件又会失去逻辑行为。两者合在一起，才是硬件设计对象。
\end{classicnote}
